\documentclass[11pt,a4paper]{article}
\usepackage[T1]{fontenc}
\usepackage{lmodern}
\usepackage[margin=23mm]{geometry}
\usepackage{amsmath,amssymb,amsthm,mathtools,booktabs,array}
\usepackage[dvipsnames]{xcolor}
\usepackage[colorlinks=true,linkcolor=MidnightBlue,urlcolor=MidnightBlue,citecolor=MidnightBlue]{hyperref}
\usepackage{microtype,fancyhdr}
\hypersetup{pdftitle={Erdos 708: Controlled Cores and Joint Divisor Certificates},pdfauthor={Research working note prepared for Denis Paliy},pdfsubject={Further partial results; the unrestricted conjectures remain unresolved}}
\pagestyle{fancy}
\fancyhf{}
\fancyhead[L]{\small Erd\H{o}s 708: controlled cores and joint certificates}
\fancyhead[R]{\small 23 November 2026}
\fancyfoot[C]{\thepage}
\setlength{\headheight}{14pt}
\setlength{\parskip}{4pt}
\setlength{\parindent}{0pt}
\newtheorem{theorem}{Theorem}[section]
\newtheorem{lemma}[theorem]{Lemma}
\newtheorem{corollary}[theorem]{Corollary}
\newtheorem{proposition}[theorem]{Proposition}
\theoremstyle{remark}
\newtheorem{remark}[theorem]{Remark}
\DeclareMathOperator{\lcm}{lcm}
\title{\textbf{Erd\H{o}s Problem 708:\\Controlled Cores and Joint Divisor Certificates}\\[5pt]
\large Further partial results and an explicit remaining inequality}
\author{Research working note prepared for Denis Paliy\\
\small AI-assisted derivations with separate internal mathematical audits}
\date{23 November 2026}
\begin{document}
\maketitle
\begin{center}
\fcolorbox{MidnightBlue}{blue!3}{\parbox{0.92\linewidth}{
\textbf{Status.} The unrestricted bounds \(2n\) and \((2+o(1))n\)
remain unproved in this investigation. The new results below close
several specific gaps in the previous argument. They include an
asymptotic coefficient-two theorem with a polylogarithmic prime core,
without a gap between prime sizes, and a separate theorem for small
and large prime bands with arbitrarily small fixed large-prime
exponent. Complete new proofs and inherited proof dependencies are
provided. Internal cross-checks are not external peer review or
formal verification; no priority claim is made.
}}
\end{center}
\begin{abstract}
We continue the investigation of product-divisibility covers of
intervals. A joint hinge comparison handles an arbitrary diffuse
core and any number of possible remainder primes, provided each
source uses at most one distinct remainder prime. A single arbitrary
core coordinate causes no error. For \(k\) substantial core
coordinates, an aggregate cap and a finite divisor expansion of
upper sets yield an error at most \(B_{\varepsilon,k}(1+|R|)\),
independent of the ambient length, with
\(B_{\varepsilon,k}\le D_\varepsilon^k\). A separate direct
construction strengthens the resulting integer source class:
\((2+o(1))n\) positions suffice when a designated core has at most
\(c_0(\log M)^2/\log\log M\) primes for a sufficiently small
absolute \(c_0>0\), and each source has at most one distinct prime
outside it. This includes every fixed power \((\log M)^a\),
\(0<a<2\). In a
different regime, a discounted signed divisor certificate allows
bounded positive overlap and absorbs arbitrary pure-core sources.
It proves the same integer bound when source primes lie below
\(\exp(o(\sqrt{\log M}))\) or above \(M^\alpha\), for every fixed
\(\alpha>0\), and for an explicit shrinking exponent. A uniform
moment calculation also excludes all sufficiently large dual means.
We isolate the remaining growing-core hinge estimate that would
imply the unrestricted asymptotic conjecture.
\end{abstract}
\begingroup
\small
\setcounter{tocdepth}{1}
\setlength{\parskip}{1pt}
\makeatletter
\renewcommand{\l@section}[2]{\@dottedtocline{1}{0em}{2em}{#1}{#2}}
\makeatother
\tableofcontents
\endgroup
\newpage

\section{Scope, common notation, and the previous checkpoint}
\label{ct:scope}

Let \(A\subseteq\{2,\ldots,M\}\) contain \(n\) distinct integers,
and let \(I=\{x+1,\ldots,x+M\}\), where \(x\ge0\) is an integer.
The smallest cardinality of a subset \(B\subseteq I\) for which
\(\prod A\mid\prod B\) is denoted by \(\tau(A,I)\).
All selected target positions are distinct. The original problem
takes \(M=\max A\); an ambient length \(M\ge\max A\) is essential
when sources are deleted during an induction.
It suffices to consider positive intervals: an interval containing
zero has a singleton cover, and an entirely negative interval is
reflected by taking absolute values.

The preceding report \cite{Synthesis} consolidated the submitted
\(32n^{5/2}\) argument, the forced-bank structural reduction, and
several fractional and integer covering theorems. The quadratic
range used here is an external theorem in Wang's written appendix:
\begin{equation}
 M\ge16n^2\quad\Longrightarrow\quad\tau(A,I)\le2n.
 \label{ct:quadratic}
\end{equation}
Its ambient formulation was checked in the preceding investigation
\cite[Appendix A]{Wang}. It is not a new result of this continuation.
Selecting the whole interval also gives \(\tau\le2n\) when
\(M\le2n\). Any remaining unrestricted exact obstruction is therefore
inside \(2n<M<16n^2\). This is a containing range, not a claim that
every smaller subrange remains untreated: Wang's result catalogue
also records a narrow large-size range just above \(2n\).

The earlier report and the inherited analytic proof notes are
included in the companion archive. In particular, precise
conditional diffuse and sieve estimates used below are stated
where needed and linked to their full inherited proofs. All new
gluing, finite-difference, overlap, transfer, and moment arguments
are proved in this report.

\subsection{A map of the additional conclusions}

\begin{center}
\small
\begin{tabular}{p{0.47\linewidth}p{0.45\linewidth}}
\toprule
Hypothesis or gap & New conclusion\\
\midrule
Diffuse core; at most one distinct outside prime per source
& One shared hinge comparison, even with arbitrarily many possible outside primes and arbitrary powers\\
One arbitrary substantial core coordinate
& The same comparison with zero additive error\\
At most \(k\) substantial core coordinates
& Error at most \(B_{\varepsilon,k}(1+|R|)\), with no dependence on \(M\)\\
A designated core of at most \(c_0(\log M)^2/\log\log M\) primes; one distinct outside prime per source
& Integer \((2+o(1))n\), without any prime-size separation\\
Arbitrary pure-core sources under the separated-support hypotheses below
& Pure-core contraction and a common hinge budget close the previous additive-excess gap\\
Prime support below \(Y\) or above \(M^\alpha\), \(\log Y=o(\sqrt{\log M})\)
& Integer \((2+o(1))n\) for every fixed \(\alpha>0\); no separate bound on large-prime multiplicity\\
Clipped dual mean at least \(241/136+o(1)\)
& Exact \(H_2\) for every interval shift\\
\bottomrule
\end{tabular}
\end{center}

\subsection{The common fractional objective}

Put \(D_p=\sum_{a\in A}v_p(a)\), and let \(P\) be the source
prime support. The raw boxed fractional cover is
\begin{equation}
 L(A,I)=\min\left\{\sum_{b\in I}z_b:
 0\le z_b\le1,\quad
 \sum_{b\in I}v_p(b)z_b\ge D_p\ (p\in P)\right\}.
 \label{ct:lp}
\end{equation}
It is feasible because
\begin{equation}
 \prod A\mid M!\mid\prod I,
 \qquad \prod I=M!\binom{x+M}{M}.
 \label{ct:supply}
\end{equation}
For weights \(w_p\ge0\), define
\[
 X_w(a)=\sum_{p\in P}w_pv_p(a),\qquad
 H_I(w)=\sum_{b\in I}(X_w(b)-1)_+,
\]
and, for a real \(c\ge0\),
\[
 \Delta_c(A,I,w)=\sum_{a\in A}(X_w(a)-c)_+-H_I(w).
\]
Ordinary linear-programming duality gives
\begin{equation}
 L(A,I)=\max_{w\ge0}
 \left\{\sum_{a\in A}X_w(a)-H_I(w)\right\}.
 \label{ct:dual}
\end{equation}
The maximum may be restricted to \(0\le w_p\le1\). Indeed, while
\(w_p>1\), every target divisible by \(p\) has load greater than one.
The slope of the displayed objective in that coordinate is
\(D_p-\sum_{b\in I}v_p(b)\le0\), by \eqref{ct:supply}.
Reducing that coordinate to one cannot decrease the objective.
Repeating this operation proves the clipping assertion.

Since \(t\le c+(t-c)_+\), a uniform estimate
\begin{equation}
 \Delta_c(A,I,w)\le E\quad(0\le w_p\le1)
 \label{ct:hinge-target}
\end{equation}
implies \(L(A,I)\le cn+E\). The proofs below compare all source
pieces against this \emph{one} target hinge. Independent covers or
independent hinge budgets cannot simply be added at coefficient two.

An optimal vertex of \eqref{ct:lp} has at most \(|P|\) coordinates
strictly between zero and one. Otherwise a nonzero perturbation
supported on those coordinates preserves the active valuation rows,
contradicting extremality. Rounding these coordinates upward gives
\begin{equation}
 \tau(A,I)\le L(A,I)+|P|.
 \label{ct:round}
\end{equation}
The fixed-ambient minimality argument below makes \(|P|=o(n)\)
where this rounding estimate is needed.

\section{A controlled core and a one-prime remainder}
\label{cg:section}

The central point of this section is a joint comparison. The core and
the remainder coordinates use one target hinge, even when many possible
remainder primes have large weights. The source restriction is that
each source uses at most one \emph{distinct} remainder prime; arbitrary
powers of that prime are permitted. No ordering or separation of prime
sizes is assumed.

All prime sets below are finite, all weights are nonnegative, and all
intervals have positive integer points. At real length $T>0$, a prefix
means $\{n\in\mathbb Z:1\le n\le T\}$ and an interval means the
integer points of $(x,x+T]$, $x\ge0$. For a prime set $F$, put
$F^\Pi=\prod_{q\in F}q$, with the empty product equal to one.

\subsection{Precise inherited analytic inputs}

We use the following conditional diffuse comparison from
\cite{Synthesis}. Its proof uses partial prime-power carriers and a
uniform interval version of the fundamental sieve lemma
\cite{Lagarias,GranvilleSound}. The statements below include the
quantifiers and constants needed here; the inherited proofs are included
in the companion archive.

\begin{proposition}[Conditional diffuse level comparison]
\label{cg:conditional-input}
There is an absolute constant $C>0$ for which the functions
\[
 e(u)=C\bigl((u/2)^{-u/4}+2^{1-u/2}\bigr)\quad(u\ge4),
 \qquad
 e_k(u)=3^ke(u/2)+3\cdot4^k2^{1-u/2}\quad(u\ge8)
\]
have the following property. Let $F$ contain at most $k$ primes, let
$L\cap F=\varnothing$, and put
\[
 V(n)=\sum_{p\in L}w_pv_p(n),\qquad
 0\le w_p\le\eta,\qquad
 \mu=\sum_{p\in L}\frac{w_p}{p-1}.
\]
For $\ell>\eta$, set
\[
 r=\lfloor\ell/\eta\rfloor\ge16,\qquad
 j=\lfloor r/2\rfloor,\qquad \kappa=r-j.
\]
If
\begin{equation}
 (1+e_k(\kappa))\mu+\eta2^{k+1-j}
 \le(\ell-\eta)(1-e_k(r)),
 \qquad e_k(r)<1,
 \label{cg:conditional-condition}
\end{equation}
then, for every $T>0$, every interval $J$ of real length $T$, and
every $s\ge0$,
\begin{equation}
 \#\{a\le T:(a,F^\Pi)=1, V(a)>s+\ell\}
 \le\#\{b\in J:(b,F^\Pi)=1, V(b)>s\}.
 \label{cg:conditional-level}
\end{equation}
The excluded primes may have arbitrary sizes, including sizes exceeding
$T$.

When $F=\varnothing$, the sharper unconditioned coefficient retained
in the same carrier proof is
\begin{equation}
 \#\{a\le T:V(a)>s+\ell\}
 \le
 \frac{(1+e(\kappa))\mu+\eta2^{1-j}}
      {(\ell-\eta)(1-e(r))}
 \#\{b\in J:V(b)>s\},
 \label{cg:unconditional-ratio}
\end{equation}
provided $e(r)<1$; here $r\ge8$ suffices.
\end{proposition}

For clarity, the numerator of the ratio in
\eqref{cg:unconditional-ratio}, divided by $\ell-\eta$, bounds
the source count per unit of a common carrier sum. The target count is
at least that carrier sum times $1-e(r)$. Thus the displayed quotient
is retained without a second loss in the comparison.

The following elementary alternative is also proved in
\cite{Synthesis}. Let $\mathbb E_0$ denote averaging over
$1,\ldots,M$. If the light support lies in $p\le M$, then
\begin{equation}
 \mu-3\eta\le\mathbb E_0V\le\mu,
 \qquad \mathbb E_0(V-\mu)^2\le8\eta\mu.
 \label{cg:moment-input}
\end{equation}
In particular, if $0<d\le1$, $\eta\le d^2/32$, and
$\mu\ge1+d$, then
\begin{equation}
 \sum_{a=1}^M\min(V(a),2)\ge M.
 \label{cg:high-mean-input}
\end{equation}
For any nonnegative raw load $X$, prime-power counts give
$\sum_I X\ge\sum_{a\le M}X(a)$. Consequently
\begin{equation}
 \sum_{a\le M}(X(a)-2)_+-\sum_{b\in I}(X(b)-1)_+
 \le M-\sum_{a\le M}\min(X(a),2).
 \label{cg:automatic}
\end{equation}
Thus \eqref{cg:high-mean-input} proves the full-prefix hinge inequality
at threshold two, and adding nonnegative coordinates preserves this
sufficient condition. Primes greater than $M$ must be removed before
computing the mean in this alternative. They change no source value;
restoring them only increases the target hinge.

\begin{lemma}[A threshold uniform in the number of exclusions]
\label{cg:uniform-eta}
For every $0<\varepsilon\le1$ there is $c_\varepsilon>0$ such that,
for every integer $k\ge0$, setting
\[
 \eta=\frac{c_\varepsilon}{k+1},\qquad \ell=1+\varepsilon
\]
makes \eqref{cg:conditional-level} valid whenever
$\mu\le1+\varepsilon/2$. We may also require
$c_\varepsilon\le\varepsilon^2/128$ for $0<\varepsilon\le1$.
\end{lemma}

\begin{proof}
For $u\ge16$,
\[
 e_k(u)\le3C3^k2^{-u/4}+6\cdot4^k2^{-u/2}.
\]
Indeed $(u/4)^{-u/8}\le2^{-u/4}$ for $u/4\ge4$.
Choose $R_\varepsilon\ge16$ so that
\[
 3C2^{-R_\varepsilon/4}+6\cdot2^{-R_\varepsilon/2}
 \le\varepsilon/16.
\]
Then $u\ge16k+R_\varepsilon$ implies
$e_k(u)\le\varepsilon/16$. Choose $c_\varepsilon$ with
\[
 c_\varepsilon\le\varepsilon^2/128,
 \qquad
 \frac{1+\varepsilon}{c_\varepsilon}
 \ge\max\{32,\,2R_\varepsilon+2\}.
\]
For $\eta=c_\varepsilon/(k+1)$, the integers in
Proposition~\ref{cg:conditional-input} satisfy
\[
 r\ge32k+2R_\varepsilon,\qquad
 \kappa\ge16k+R_\varepsilon,\qquad j\ge k+1.
\]
Hence both sieve errors are at most $\varepsilon/16$, and
$\eta2^{k+1-j}\le\eta$. At $\mu\le1+\varepsilon/2$, the
right-hand lower bound in \eqref{cg:conditional-condition} minus
its left-hand upper bound is at least
\[
 (1+\varepsilon-\eta)(1-\varepsilon/16)
 -(1+\varepsilon/16)(1+\varepsilon/2)-\eta
 \ge\frac{46}{128}\varepsilon-\frac{191}{2048}\varepsilon^2>0.
\]
Here $\eta\le\varepsilon/128$ was used. This proves the uniform
choice. Subsequent shrinking of $c_\varepsilon$ is harmless.
\end{proof}

\subsection{Gluing increments into one target hinge}

\begin{lemma}[One-prime source gluing]
\label{cg:gluing}
Let $R$ be a finite prime set, let $C(n)\ge0$ satisfy
$C(r^jt)=C(t)$ for every $r\in R$ and $j\ge0$, and let
$u_r:\mathbb Z_{\ge0}\to[0,1]$ be nondecreasing with $u_r(0)=0$.
Put $U(n)=\sum_{r\in R}u_r(v_r(n))$, $\ell>0$, and $c=1+\ell$.
Suppose $0<\rho\le1$ and $A_0,B_0\ge0$ have these properties:
\begin{align}
 \sum_{a\le M}(C(a)-c)_+
 &\le\rho\sum_{b\in I}(C(b)-1)_++A_0,
 \label{cg:base-assumption}\\
 \#\{a\le N:C(a)>s+\ell\}
 &\le\rho\#\{b\in J:C(b)>s\}+B_0
 \label{cg:level-assumption}
\end{align}
for all positive integer $N$, all $N$-term intervals $J$, and all
$0\le s\le1$. If $A\subseteq[2,M]$ is distinct and each source
has at most one distinct $R$-prime, then
\begin{equation}
 \sum_{a\in A}(C(a)+U(a)-c)_+
 \le\rho\sum_{b\in I}(C(b)+U(b)-1)_+
       +A_0+B_0|R|.
 \label{cg:gluing-conclusion}
\end{equation}
More precisely, the last error may be replaced by
$A_0+\sum_{r\in R}\min\{B_0,m_r\}$, where
$m_r=\#\{a\in A:r\mid a\}$.
\end{lemma}

\begin{proof}
For $h\in\{c,1\}$, define the nonnegative increment
\[
 D_{h,r}(n)=(C(n)+u_r(v_r(n))-h)_+-(C(n)-h)_+.
\]
Telescoping the profile and integrating the derivative of the hinge
gives the exact identity
\begin{equation}
 D_{h,r}(n)=\sum_{j\ge1}\mathbf1_{r^j\mid n}
  \int_{u_r(j-1)}^{u_r(j)}\mathbf1_{C(n)>h-v}\,dv.
 \label{cg:increment}
\end{equation}
For each $j$, divide multiples of $r^j$ by $r^j$. The source
quotients form a prefix of length $N=\lfloor M/r^j\rfloor$,
whereas the target quotients form an interval with at least $N$
consecutive positions. When $N>0$, take an $N$-term subinterval
and apply \eqref{cg:level-assumption} at $s=1-v\in[0,1]$.
The unused target positions contribute nonnegatively. When $N=0$,
the source count is zero. Invariance of $C$ under $r$-powers and
\eqref{cg:increment} now give
\begin{equation}
 \sum_{a\le M}D_{c,r}(a)
 \le\rho\sum_{b\in I}D_{1,r}(b)+B_0,
 \label{cg:per-prime}
\end{equation}
because $\sum_{j\ge1}(u_r(j)-u_r(j-1))\le1$.
Thus the error is paid once per prime, not once per prime power.

At each source, at most one remainder coordinate is nonzero, so
\begin{equation}
 (C(a)+U(a)-c)_+=(C(a)-c)_++\sum_{r\in R}D_{c,r}(a).
 \label{cg:source-identity}
\end{equation}
This identity is only needed on the source set. Every summand is
nonnegative, so distinctness permits us to enlarge source sums to the
prefix when applying \eqref{cg:base-assumption} and
\eqref{cg:per-prime}.

For $x,z_i\ge0$, one has
\begin{equation}
 (x-1)_++\sum_i\bigl((x+z_i-1)_+-(x-1)_+\bigr)
 \le\left(x+\sum_i z_i-1\right)_+.
 \label{cg:superadditive}
\end{equation}
If $x\ge1$ this is equality. Otherwise put $t=1-x>0$ and use
$\sum_i(z_i-t)_+\le(\sum_i z_i-t)_+$. Applying
\eqref{cg:superadditive} pointwise on the targets proves
\eqref{cg:gluing-conclusion}; the base target hinge has been used once.

For the refinement, the actual source increment
$S_r=\sum_{a\in A}D_{c,r}(a)$ is at most $m_r$, since each
increment is at most one and vanishes off the $r$-multiples.
It is also at most $\rho\sum_I D_{1,r}+B_0$ by
\eqref{cg:per-prime}. These two bounds imply
$S_r\le\rho\sum_I D_{1,r}+\min\{B_0,m_r\}$, proving the
sharper assertion.
\end{proof}

\begin{lemma}[Simultaneous restoration of raw tails]
\label{cg:tails}
Let $P_0$ be any finite set of coordinates capped at one, and write
\[
 Y(n)=V(n)+\sum_{p\in P_0}\min\{w_pv_p(n),1\},\qquad
 T(n)=\sum_{p\in P_0}(w_pv_p(n)-1)_+,
\]
where $V(n)\ge0$. If, for a distinct prefix subset $A$,
\[
 \sum_{a\in A}(Y(a)-c)_+\le\sum_{b\in I}(Y(b)-1)_++E_0,
\]
then the same estimate holds with $Y$ replaced by $Y+T$.
\end{lemma}

\begin{proof}
Pointwise, $(Y+T-c)_+\le T+(Y-c)_+$. Moreover,
\[
 (Y+T-1)_+=T+(Y-1)_+,
\]
because $T>0$ forces at least one capped coordinate to equal one,
and hence $Y\ge1$. Each summand of $T$ is a nondecreasing
prime-power profile vanishing at exponent zero. Expanding it in
nonnegative increments times $\mathbf1_{p^j\mid n}$ shows
\[
 \sum_{a\in A}T(a)\le\sum_{a\le M}T(a)\le\sum_{b\in I}T(b).
\]
Combining these inequalities proves the assertion. All tails are
restored in this single comparison.
\end{proof}

\subsection{Diffuse cores and one unrestricted core coordinate}

\begin{theorem}[A diffuse core with arbitrary one-prime remainders]
\label{thm:cg-diffuse}
For every $\varepsilon>0$ there is $\eta_\varepsilon>0$ with the
following property. Let $E,R$ be disjoint finite prime sets and put
\[
 X(n)=V(n)+W(n),\qquad
 V(n)=\sum_{p\in E}w_pv_p(n),\qquad
 W(n)=\sum_{r\in R}w_rv_r(n),
\]
where $0\le w_p\le\eta_\varepsilon$ for $p\in E$ and
$w_r\ge0$ is unrestricted for $r\in R$. If
$A\subseteq\{2,\ldots,M\}$ is distinct and each source is
divisible by at most one distinct $R$-prime, then, for every shift,
\begin{equation}
 \sum_{a\in A}(X(a)-2-\varepsilon)_+
 \le\sum_{b\in I}(X(b)-1)_+.
 \label{cg:diffuse-hinge}
\end{equation}
Arbitrary powers of that one remainder prime and arbitrary unweighted
cofactors are allowed. All thresholds are independent of the sizes and
cardinalities of the prime sets.
\end{theorem}

\begin{proof}
It suffices to treat $0<\varepsilon\le1$. Put
$\ell=1+\varepsilon$, $c=2+\varepsilon$, and
\[
 \rho_\varepsilon=\frac{1+3\varepsilon/4}{1+\varepsilon}<1.
\]
Choose $\eta=\eta_\varepsilon\le\varepsilon^2/128$ so small
that, with $r,j,\kappa$ as in
Proposition~\ref{cg:conditional-input},
\begin{equation}
 \frac{(1+e(\kappa))(1+\varepsilon/2)+\eta2^{1-j}}
      {(\ell-\eta)(1-e(r))}\le\rho_\varepsilon,
 \qquad e(r)<1,\quad r\ge8.
 \label{cg:strict-choice}
\end{equation}
Such a choice exists: the ratio tends to
$(1+\varepsilon/2)/(1+\varepsilon)$ as $\eta\downarrow0$.

Discard core primes greater than $M$, and compute
$\mu=\sum_{p\in E}w_p/(p-1)$ on the remaining support.
If $\mu\ge1+\varepsilon/2$, equations
\eqref{cg:high-mean-input} and \eqref{cg:automatic} give the
stronger full-prefix threshold-two inequality after adding $W$.
Restriction to $A$ proves the assertion.

Suppose $\mu\le1+\varepsilon/2$. Equations
\eqref{cg:unconditional-ratio} and \eqref{cg:strict-choice} give
\[
 \#\{a\le T:V(a)>s+\ell\}
 \le\rho_\varepsilon\#\{b\in J:V(b)>s\}
 \qquad(T>0, s\ge0).
\]
This comparison is uniform at every smaller length as well.
Integrating over $s\ge1$ gives its base hinge version at thresholds
$c$ and one. Apply Lemma~\ref{cg:gluing} with $C=V$,
$u_r(e)=\min\{w_re,1\}$, and $A_0=B_0=0$. This proves
\begin{equation}
 \sum_{a\in A}(V(a)+U(a)-c)_+
 \le\rho_\varepsilon\sum_{b\in I}(V(b)+U(b)-1)_+.
 \label{cg:strict-capped}
\end{equation}
Dropping the factor $\rho_\varepsilon\le1$ and applying
Lemma~\ref{cg:tails} to all $R$-coordinates restores their arbitrary
raw weights and powers. Finally restore the removed core primes on the
target side. This proves \eqref{cg:diffuse-hinge}.
\end{proof}

In the low-mean case, if every source has at most one remainder
\emph{occurrence} and $w_r\le1$, its remainder tail is zero.
Then \eqref{cg:strict-capped} yields the strict raw comparison with
factor $\rho_\varepsilon$ as well. Arbitrary remainder powers need
only the non-strict comparison in the theorem.

\begin{theorem}[One arbitrary core coordinate]
\label{thm:cg-one-heavy}
For every $\varepsilon>0$ there is $\eta_\varepsilon^*>0$ such
that Theorem~\ref{thm:cg-diffuse} remains valid when one prime
$q\in E$ has an arbitrary nonnegative raw weight and every other
core weight is at most $\eta_\varepsilon^*$. The number of
possible remainder primes and their weights are still unrestricted.
\end{theorem}

\begin{proof}
Again take $0<\varepsilon\le1$ and $\ell=1+\varepsilon$.
Write
\[
 V_0(n)=\sum_{p\in E\setminus\{q\}}w_pv_p(n),\qquad
 C(n)=V_0(n)+w_qv_q(n).
\]
Choose $\eta_\varepsilon^*$ no larger than the $k=1$ threshold
in Lemma~\ref{cg:uniform-eta}. Discard light primes above $M$
before computing their mean $\mu_0$.
The high-mean case $\mu_0\ge1+\varepsilon/2$ follows from
\eqref{cg:high-mean-input} and \eqref{cg:automatic} after adding
both the $q$-coordinate and the remainder load.

In the low-mean case, the conditional input gives
\begin{equation}
 \#\{a\le T:(a,q)=1, V_0(a)>s+\ell\}
 \le\#\{b\in J:(b,q)=1, V_0(b)>s\}
 \quad(s\ge0).
 \label{cg:one-conditional}
\end{equation}
We claim that the unrestricted core $C$ itself satisfies
\begin{equation}
 \#\{a\le T:C(a)>s+\ell\}
 \le\#\{b\in J:C(b)>s\}
 \quad(T>0, s\ge0).
 \label{cg:one-level}
\end{equation}
Decompose integers uniquely as $q^jt$ with $(t,q)=1$.
For every $j$ with $s-w_qj\ge0$, apply
\eqref{cg:one-conditional} at real length $T/q^j$, to the scaled
target interval and shifted level $s-w_qj$.
If $w_q=0$, these are all pieces. Otherwise the negative shifted
levels form the single suffix
\[
 j\ge j_0,\qquad j_0=\lfloor s/w_q\rfloor+1.
\]
On this suffix the target condition is automatic, since $V_0\ge0$.
The source count is at most $\lfloor T/q^{j_0}\rfloor$, whereas
the target count of the entire suffix is exactly the number of
$q^{j_0}$-multiples in $J$, which is at least that number.
All target exponents belong to this comparison, including $q^j>T$.
The boundary $w_qj=s$ belongs to the nonnegative-level case, not
to the suffix. Summing proves \eqref{cg:one-level} without error.

Integrating \eqref{cg:one-level} over $s\ge1$ gives the base
hinge comparison. The function $C$ is nonnegative and invariant
under every remainder-prime power. Lemma~\ref{cg:gluing}, now
with $\rho=1$ and zero errors, glues all capped remainder profiles.
Lemma~\ref{cg:tails} restores their tails. This proves the result.
The same suffix argument would permit any nondecreasing nonnegative
$q$-adic profile, although only the raw profile is needed here.
\end{proof}

\subsection{Finite variation of a capped exceptional upper set}

For several exceptional core coordinates, the negative-level pieces
need not form a single prime-power suffix. Their collective discrepancy
can nevertheless be bounded independently of the interval length.

\begin{lemma}[A finite divisor expansion]
\label{lem:cg-variation}
Let $F$ contain at most $k$ primes, each with $w_q\ge\eta>0$, and
put
\[
 Z(n)=\sum_{q\in F}\min\{w_qv_q(n),1\},\qquad
 f_s(n)=\mathbf1_{Z(n)>s}\quad(s\ge0).
\]
For every real $T>0$ and interval $J$ of length $T$,
\begin{equation}
 \left|\sum_{a\le T}f_s(a)-\sum_{b\in J}f_s(b)\right|
 \le4^k\binom{k+\lfloor s/\eta\rfloor}{k}.
 \label{cg:variation-bound}
\end{equation}
If $s\ge k$, or if $F$ is empty, the difference is zero.
\end{lemma}

\begin{proof}
The empty case is immediate. Write $F=\{q_1,\ldots,q_f\}$,
$f\le k$, and set $b_i=\lceil1/w_{q_i}\rceil$.
Let $u_i(j)=\min\{w_{q_i}j,1\}$ and
\[
 g_s(j_1,\ldots,j_f)=\mathbf1_{\sum_i u_i(j_i)>s}.
\]
For a nonnegative integer vector $\mathbf j$, define its mixed
divisor coefficient by
\begin{equation}
 a_{\mathbf j}(s)=
 \sum_{\substack{e_i\in\{0,1\}\\ e_i\le j_i\ (1\le i\le f)}}
 (-1)^{\sum_i e_i}g_s(\mathbf j-\mathbf e).
 \label{cg:finite-difference}
\end{equation}
Thus subtraction of one is allowed only in positive coordinates.
Telescoping separately in every coordinate gives the exact identity
\begin{equation}
 f_s(n)=\sum_{\mathbf j:\,\prod_iq_i^{j_i}\mid n}a_{\mathbf j}(s).
 \label{cg:divisor-expansion}
\end{equation}
If $j_i>b_i$ for any coordinate, the two adjacent layers of that
capped profile coincide, and its mixed difference is zero.
Consequently the expansion is finite. Also
$|a_{\mathbf j}(s)|\le2^f$.

For a nonzero vector $\mathbf j$, put $z_i=(j_i-1)_+$.
If $\sum_i u_i(z_i)>s$, all corners in
\eqref{cg:finite-difference} equal one, and their alternating sum
vanishes. The zero vector also has coefficient zero, since $s\ge0$
and $g_s(\mathbf0)=0$. A nonzero coefficient therefore requires
\[
 \sum_i u_i(z_i)\le s.
\]
On the relevant grid $j_i\le b_i$, one has $z_i<b_i$, so
$u_i(z_i)=w_{q_i}z_i\ge\eta z_i$. Hence every nonzero
coefficient satisfies
\begin{equation}
 \sum_i z_i\le s/\eta.
 \label{cg:corner-support}
\end{equation}
There are $\binom{f+\lfloor s/\eta\rfloor}{f}$ possible vectors
of nonnegative integers $z_i$ satisfying \eqref{cg:corner-support}.
Each zero coordinate $z_i=0$ permits $j_i=0$ or $1$; a positive
$z_i$ determines $j_i$ uniquely. Thus
\begin{equation}
 \sum_{\mathbf j}|a_{\mathbf j}(s)|
 \le4^f\binom{f+\lfloor s/\eta\rfloor}{f}
 \le4^k\binom{k+\lfloor s/\eta\rfloor}{k}.
 \label{cg:variation-norm}
\end{equation}

For each divisor $d$, its count in $J$ is either
$\lfloor T/d\rfloor$ or $\lfloor T/d\rfloor+1$.
Its prefix count is $\lfloor T/d\rfloor$. Each divisor-count
discrepancy thus has absolute value at most one. The triangle
inequality and \eqref{cg:variation-norm} prove
\eqref{cg:variation-bound} in both signs. This two-sided version
will be needed when the aggregate core tail is restored.
Finally, $Z(n)\le f\le k$, so $f_s$ vanishes for $s\ge k$.
\end{proof}

\subsection{Aggregating negative levels before estimating them}

\begin{lemma}[The collective exceptional-level defect]
\label{lem:cg-level-defect}
Let \(F,L\) be disjoint, let \(F\) contain at most \(k\) primes with
\(w_q\ge\eta>0\), and define
\[
 \begin{aligned}
 V(n)&=\sum_{p\in L}w_pv_p(n),&
 X_F(n)&=\sum_{q\in F}w_qv_q(n),\\
 Z(n)&=\min\{X_F(n),1\},& C(n)&=V(n)+Z(n).
 \end{aligned}
\]
Suppose the conditional comparison \eqref{cg:conditional-level}
holds for \(V,F,\ell\) at every nonnegative level and every real
length. Then, for all \(s\ge0\), all real \(T>0\), and intervals
\(J\) of length \(T\),
\begin{equation}
 \begin{split}
 &\#\{a\le T:C(a)>s+\ell\}-\#\{b\in J:C(b)>s\}\\
 &\qquad\le
 \sum_{a\le T}\mathbf1_{Z(a)>s}
       -\sum_{b\in J}\mathbf1_{Z(b)>s}.
 \end{split}
 \label{cg:aggregate-defect}
\end{equation}
In particular, the level defect for \(0\le s\le1\) is at most
\[
 B_{\eta,k}=4^k\binom{k+\lfloor1/\eta\rfloor}{k}
 \quad(k\ge1),
\]
and, with \(c=1+\ell\), the base hinge has no defect:
\begin{equation}
 \sum_{a\le M}(C(a)-c)_+-\sum_{b\in I}(C(b)-1)_+
 \le0.
 \label{cg:base-defect}
\end{equation}
The aggregate cap also satisfies
\begin{equation}
 \left|\sum_{a\le T}Z(a)-\sum_{b\in J}Z(b)\right|
 \le B_{\eta,k}.
 \label{cg:aggregate-cap-variation}
\end{equation}
For \(k=0\), \(B_{\eta,0}=0\).
\end{lemma}

\begin{proof}
Decompose each positive integer uniquely as \(n=ht\), where \(h\)
is its full \(F\)-smooth part and \((t,F^\Pi)=1\). Thus
\[
 C(ht)=Z(h)+V(t).
\]
For every \(h\) with \(Z(h)\le s\), apply the conditional input at
length \(T/h\), in the scaled target interval, at the nonnegative
level \(s-Z(h)\). The defect from this exact part is nonpositive.

For \(h\) with \(Z(h)>s\), the target condition holds automatically,
because \(V(t)\ge0\). The source contribution is at most the
unrestricted prefix count with this exact \(F\)-part. Sum all these
pieces before estimating their discrepancy. The sum is precisely
the difference of the two exceptional upper-set counts on the
right of \eqref{cg:aggregate-defect}. All target exact parts are
included, even those with \(h>T\). This proves the displayed
comparison; no error has been assigned separately to each smooth
part.

For \(0\le s<1\), the exceptional upper set is unchanged by
individual coordinate caps:
\[
 \mathbf1_{Z(n)>s}
 =\mathbf1_{X_F(n)>s}
 =\mathbf1_{\sum_{q\in F}\min\{w_qv_q(n),1\}>s}.
\]
Lemma~\ref{lem:cg-variation} therefore bounds its discrepancy in
both signs by \(B_{\eta,k}\). At \(s\ge1\), the upper set is
empty. This proves the level error on \(s\in[0,1]\).
Integrating \eqref{cg:aggregate-defect} over \(s\ge1\) gives
\eqref{cg:base-defect}, since its right side is now identically
zero. Finally,
\[
 Z(n)=\int_0^1\mathbf1_{Z(n)>s}\,ds.
\]
Integrating the two-sided discrepancy bound over this unit interval
proves \eqref{cg:aggregate-cap-variation}.
\end{proof}

\subsection{The uniform controlled-core theorem}

\begin{theorem}[Controlled substantial core coordinates]
\label{thm:cg-controlled}
For every \(\varepsilon>0\) there is \(c_\varepsilon>0\) such
that the following holds for every integer \(k\ge0\). Put
\[
 \eta=\frac{c_\varepsilon}{k+1}.
\]
Let \(E,R\) be disjoint finite prime sets. All but at most \(k\)
weights in \(E\) are at most \(\eta\); the remaining core weights
and all \(R\)-weights may be arbitrary nonnegative numbers. Define
\[
 X(n)=\sum_{p\in E\cup R}w_pv_p(n).
\]
If \(A\subseteq\{2,\ldots,M\}\) is distinct and each source is
divisible by at most one distinct \(R\)-prime, then
\begin{equation}
 \sum_{a\in A}(X(a)-2-\varepsilon)_+
       -\sum_{b\in I}(X(b)-1)_+
 \le B_{\varepsilon,k}(1+|R|)
 \label{cg:controlled-defect}
\end{equation}
for every positive \(M\)-term interval. Arbitrary powers of the
one remainder prime and arbitrary unweighted cofactors are allowed.
For \(k\ge1\), one may take
\begin{equation}
 B_{\varepsilon,k}
 =4^k\binom{k+\lfloor1/\eta\rfloor}{k}.
 \label{cg:error-constants}
\end{equation}
For \(k=0\), take \(B_{\varepsilon,0}=0\). The constant is
independent of \(M\), the prime values, the weights, and the interval
shift. One may restrict \(R\) to the primes occurring in the
sources. More precisely, the whole right side of
\eqref{cg:controlled-defect} may be replaced by
\begin{equation}
 B_{\varepsilon,k}+\sum_{r\in R}\min\{B_{\varepsilon,k},m_r\},
 \qquad m_r=\#\{a\in A:r\mid a\}.
 \label{cg:fibre-refinement}
\end{equation}
\end{theorem}

\begin{proof}
It suffices to treat \(0<\varepsilon\le1\). Choose
\(c_\varepsilon\) as in Lemma~\ref{cg:uniform-eta} and put
\[
 F=\{p\in E:w_p>\eta\},\qquad L=E\setminus F,\qquad
 V(n)=\sum_{p\in L}w_pv_p(n).
\]
Then \(|F|\le k\), and every \(F\)-weight is at least \(\eta\).
Discard light primes above \(M\) before computing
\(\mu_L=\sum_{p\in L}w_p/(p-1)\). If
\(\mu_L\ge1+\varepsilon/2\), the moment alternative
\eqref{cg:high-mean-input} applies because
\(\eta\le\varepsilon^2/128\). Adding all the remaining
nonnegative coordinates and using \eqref{cg:automatic} proves
\eqref{cg:controlled-defect} with zero error.

It remains to assume \(\mu_L\le1+\varepsilon/2\). The conditional
level comparison is now available uniformly at every real length,
with excluded set \(F\) and gap \(\ell=1+\varepsilon\).
Cap the whole substantial core at one, and cap each remainder
coordinate at one:
\[
 \begin{split}
 X_F(n)&=\sum_{q\in F}w_qv_q(n),\qquad
 Z(n)=\min\{X_F(n),1\},\qquad T_F(n)=X_F(n)-Z(n),\\
 U(n)&=\sum_{r\in R}\min\{w_rv_r(n),1\},\\
 T_R(n)&=\sum_{r\in R}(w_rv_r(n)-1)_+,\qquad
 C(n)=V(n)+Z(n).
 \end{split}
\]
Thus \(X=C+U+T_F+T_R\). Lemma~\ref{lem:cg-level-defect}
gives a level error \(B_{\varepsilon,k}\) on \(s\in[0,1]\)
and zero base hinge error. The function \(C\) is
invariant under every \(R\)-prime power. Lemma~\ref{cg:gluing},
with \(\rho=1\) and \(A_0=0\), therefore gives
\[
 \sum_{a\in A}(C(a)+U(a)-2-\varepsilon)_+
 \le\sum_{b\in I}(C(b)+U(b)-1)_+
       +B_{\varepsilon,k}|R|.
\]

The aggregate core tail is not a sum of separate prime-power
profiles, so it requires a separate justified estimate. Linear raw
prime-power counting and \eqref{cg:aggregate-cap-variation} give
\begin{equation}
 \begin{split}
 \sum_{a\le M}T_F(a)-\sum_{b\in I}T_F(b)
 &=\sum_{a\le M}X_F(a)-\sum_{b\in I}X_F(b)
       +\sum_{b\in I}Z(b)-\sum_{a\le M}Z(a)\\
 &\le B_{\varepsilon,k}.
 \end{split}
 \label{cg:aggregate-tail}
\end{equation}
Since \(T_F\ge0\), restriction to \(A\) preserves this upper
bound. The remainder tails retain exact domination,
\(\sum_A T_R\le\sum_I T_R\), as in Lemma~\ref{cg:tails}.
If \(T_F+T_R>0\), then \(C+U\ge1\): either \(Z=1\), or
one capped remainder coordinate equals one. Consequently
\[
 (X-1)_+=(T_F+T_R)+(C+U-1)_+,
 \qquad
 (X-2-\varepsilon)_+
 \le T_F+T_R+(C+U-2-\varepsilon)_+.
\]
Restoring these tails therefore costs at most one additional
\(B_{\varepsilon,k}\), not a separate cost for each exceptional
coordinate or smooth part. Restoring light primes above \(M\)
only increases the target load. This proves
\eqref{cg:controlled-defect}.

Primes of \(R\) absent from \(A\) may be removed before the proof
and restored only on the target side. The refined error
\eqref{cg:fibre-refinement} follows from the refined assertion of
Lemma~\ref{cg:gluing}; the aggregate tail restoration adds the
single \(B_{\varepsilon,k}\) in \eqref{cg:fibre-refinement}.
\end{proof}

The coefficient in the theorem is not asserted to be optimal.
In particular, Theorem~\ref{thm:cg-one-heavy} proves zero error
when \(k=1\) and the light threshold is chosen for that theorem.
The finite-variation estimate is designed to remain uniform for a
growing number of substantial core coordinates.

The error in \eqref{cg:aggregate-tail} cannot simply be replaced by
an assertion of exact prefix dominance for aggregate tails. Take
\(F=\{11,13,17\}\), all weights one, \(M=221\), and
\(I=\{2313,\ldots,2533\}\). The tail
\(T_F=(X_F-1)_+\) has prefix sum five: its positive contributions
occur at \(121,143,169,187,221\), each with value one. On \(I\),
the sum is four: \(2366\) and \(2420\) contribute one each, while
\(2431=11\cdot13\cdot17\) contributes two. This is a counterexample
to that intermediate dominance assertion, not to a hinge theorem
or to problem 708.

\begin{corollary}[Exponential growth of the error constant]
\label{cg:error-growth}
For \(k\ge1\), the constant above satisfies
\begin{equation}
 B_{\varepsilon,k}\le D_\varepsilon^k,\qquad
 D_\varepsilon=4e(1+2/c_\varepsilon).
 \label{cg:quantitative-errors}
\end{equation}
In the notation \(A_{\varepsilon,k}+B_{\varepsilon,k}|R|\),
one may take \(A_{\varepsilon,k}=B_{\varepsilon,k}\).
In particular, with no remainder primes, at most \(k\) substantial
weights cause a full-prefix hinge defect at most \(D_\varepsilon^k\),
uniformly in the ambient length.
\end{corollary}

\begin{proof}
For \(k\ge1\),
\[
 \lfloor1/\eta\rfloor\le2k/c_\varepsilon.
\]
Apply \(\binom Nk\le(eN/k)^k\) in
\eqref{cg:error-constants}. The no-remainder assertion follows
by taking the full distinct prefix as the source set in the theorem.
\end{proof}

\subsection{What the controlled estimate excludes}

Let \(n=|A|\to\infty\), and suppose the number of active remainder
primes satisfies \(|R|\le C_0n/\log n\), for a fixed constant
\(C_0\). Fix \(\varepsilon>0\) and \(0<\gamma<1\). If
\begin{equation}
 k\le\frac{(1-\gamma)\log\log n}{\log D_\varepsilon},
 \qquad
 \#\{p\in E:w_p>c_\varepsilon/(k+1)\}\le k,
 \label{cg:controlled-range}
\end{equation}
then
\[
 B_{\varepsilon,k}|R|
 =O\!\left(\frac{n}{(\log n)^\gamma}\right)=o(n),
 \qquad B_{\varepsilon,k}=O((\log n)^{1-\gamma})=o(n).
\]
Since \(x\le2+\varepsilon+(x-2-\varepsilon)_+\), the raw
dual objective of every such vector is at most
\begin{equation}
 \sum_{a\in A}X(a)-\sum_{b\in I}(X(b)-1)_+
 \le(2+\varepsilon)n+o(n).
 \label{cg:dual-exclusion}
\end{equation}
All error bounds here are independent of \(M\); no relation
between \(M\) and \(n\) is needed for this deduction. In the
minimal-obstruction class described later, the prime bound
\(p<cn\), with fixed \(c>2\), gives the required bound on
\(|R|\) by Chebyshev's estimate.

In particular, \eqref{cg:controlled-range} excludes a witness with
objective at least \((2+2\varepsilon)n\) for all sufficiently large
\(n\). It does not exclude an arbitrarily small excess over
\((2+\varepsilon)n\). This distinction is important when interpreting
an asymptotic estimate as a restriction on possible witnesses.
The theorem permits arbitrarily many substantial remainder weights;
the controlled count in \eqref{cg:controlled-range} is entirely
inside the designated core.

There is also an integer source-class consequence, proved with the
rounding step in Corollary~\ref{cor:tr-small-core}. If a designated
core has size \(o(\log\log M)\) and each source has at most one
distinct prime outside it, then
\(\tau(A,I)\le(2+o(1))|A|\), uniformly, without any separation of
prime sizes. In that application all core weights are unrestricted:
the number of substantial ones is bounded by the size of the core
for every dual vector.

That integer corollary records what follows from the controlled-weight
method. The separate construction in
Theorem~\ref{thm:sc-polylog} permits a designated core of size
\((\log M)^a\) for any fixed \(a<2\), and
Corollary~\ref{cor:sc-chebyshev-core} permits
\(c_0(\log M)^2/\log\log M\) core primes for an absolute \(c_0>0\).
Those constructions do not, by themselves,
increase the number of substantial coordinates handled uniformly
inside an unrestricted larger core in \eqref{cg:controlled-range}.

The unrestricted demand core can have size
\(O(\sqrt{n/\log n})\), much larger than the controlled range above.
The finite-variation errors therefore do not establish the
unrestricted asymptotic conjecture. A general proof would have to
bound their joint contribution more economically, or bypass these
upper-set discrepancies while retaining the single target budget.

\section{Integer transfer at a fixed ambient length}
\label{sec:tr-transfer}

The source-specific hinge estimates above concern fractional covers. This
section gives the integer step with the ambient length retained under
deletion. No generic assertion that fractional covers have negligible
rounding cost is needed: that cost is bounded at a hypothetical minimal
failing subset by forcing its source primes to be small.

For $A\subseteq\{2,\ldots,M\}$ and
$I=\{x+1,\ldots,x+M\}$, $x\in\mathbb Z_{\ge0}$, write
\[
 D_p(A)=\sum_{a\in A}v_p(a).
\]
The raw fractional cost is
\[
 \mathcal L_{\rm raw}(A,I)=
 \min\left\{\sum_{b\in I}y_b:
 0\le y_b\le1,\quad
 \sum_{b\in I}v_p(b)y_b\ge D_p(A)\ (p\in P)
 \right\}.
\]
Here $P$ is the source prime support. The LP is feasible because
$\prod A\mid M!\mid\prod I$. We use the
external quadratic theorem of \cite{Wang} in its ambient form:
\begin{equation}
 M\ge16|A|^2\quad\Longrightarrow\quad\tau(A,I)\le2|A|.
 \label{eq:tr-quadratic}
\end{equation}
The written proof uses only $a\le M$, so equality $M=\max A$ is not
required. If an optimal basic feasible boxed solution has $r$ covering
rows, at most $r$ coordinates are fractional. Rounding them up preserves
feasibility and increases the objective by at most $r$.

\begin{lemma}[Hereditary fractional-to-integer transfer]
\label{lem:tr-hereditary}
For each $M$, let $\mathcal C(M)$ be a family of subsets of
$\{2,\ldots,M\}$, closed under taking subsets while retaining $M$.
Fix $\delta>0$. Suppose that, for every sufficiently large ambient $M$,
uniformly over $A\in\mathcal C(M)$ and every shift,
\[
 \mathcal L_{\rm raw}(A,I)\le(2+\delta)|A|.
\]
Then, for every sufficiently large ambient $M$, uniformly over the same
class and shifts,
\[
 \tau(A,I)\le(2+2\delta)|A|.
\]
If the fractional hypothesis holds for every $\delta>0$, the integer
conclusion is $\tau(A,I)\le(2+o(1))|A|$ as $|A|\to\infty$.
\end{lemma}

\begin{proof}
Put $c=2+2\delta$. Suppose a source set fails the bound $c|A|$, and
choose a failing subset $B$ of least positive cardinality $m$, in the
same interval. Each deletion has a cover of size at most
\[
 \lfloor c(m-1)\rfloor=K-1.
\]
If a source $a\in B$ has a prime divisor $p\ge K$, its $a/p$ bank
contains at least $\lfloor M/(a/p)\rfloor\ge p\ge K$ positions.
One is unused by the deletion cover. Add it and, if needed, one further
unused position supplying the missing $p$-unit. Such a repair exists by
full-interval supply. The result has at most
\[
 \lfloor c(m-1)\rfloor+2\le\lfloor cm\rfloor
\]
positions, a contradiction. Hence all active source primes are below
$K\le cm$, and there are at most $\pi(cm)$ rows in the raw LP.

Heredity permits the fractional hypothesis to be applied to $B$ at the
original $M$. Basic-solution rounding gives
\begin{equation}
 \tau(B,I)\le(2+\delta)m+\pi(cm).
 \label{eq:tr-basic-round}
\end{equation}
Choose $m_0$ such that $\pi(cm)\le\delta m$ for every $m\ge m_0$;
Chebyshev's estimate suffices. A failing subset cannot satisfy
$M\ge16m^2$, by \eqref{eq:tr-quadratic}. Therefore
\begin{equation}
 m>\sqrt M/4.
 \label{eq:tr-minimal-large}
\end{equation}
Once $M$ exceeds both the fractional threshold and $16m_0^2$,
\eqref{eq:tr-basic-round} contradicts failure. This proves the lemma.
Finally, $|A|\to\infty$ implies $M\ge|A|+1\to\infty$, giving the
stated uniform asymptotic conclusion.
\end{proof}

No additive offset is used in minimality. An all-cardinality bound
$(2+\varepsilon)|A|+C_\varepsilon$ may be obtained afterward: below the
large-$M$ threshold, the whole interval is a cover of size at most $M$.
This order avoids any circular dependence of that threshold on an offset.

\subsection{An arbitrary small core with one outside prime per source}

The additive error in Theorem~\ref{thm:cg-controlled} can also be handled
directly inside the same minimality argument. For fixed $\delta>0$, denote
its sharp constant by $B_{\delta,k}$, with
\begin{equation}
 B_{\delta,k}\le D_\delta^k,
 \label{eq:tr-error-growth}
\end{equation}
and $B_{\delta,0}=0$. The hinge defect is at most
$B_{\delta,k}(1+|R|)$; the first term accounts for the aggregate core tail.

\begin{corollary}[A core of size $o(\log\log M)$]
\label{cor:tr-small-core}
Let $k(M)$ be a nonnegative integer-valued function with
$k(M)=o(\log\log M)$ as $M\to\infty$. Suppose that a distinct source
set $A\subseteq\{2,\ldots,M\}$ admits a set $E$ of at most $k(M)$
primes such that each source is divisible by at most one distinct prime
outside $E$. Arbitrary powers of that outside prime are allowed.
Then, uniformly over all such sets, actual prime values, ambient lengths,
and shifts,
\[
 \tau(A,I)\le(2+o(1))|A|\qquad(|A|\to\infty).
\]
No separation between core and outside prime sizes is required.
\end{corollary}

\begin{proof}
Fix $\delta>0$ and put $c=2+2\delta$. Suppose that $A$ fails the bound
$c|A|$. Choose a failing subset $B$ of least cardinality $m$, retaining
the original $M$ and the same core $E$. The one-distinct-outside-prime
property is hereditary with these parameters fixed. The preceding prime
forcing argument shows that every active source prime is below $cm$.
Let $R_B$ consist only of the outside primes actually occurring in $B$.
Then
\begin{equation}
 |R_B|\le\pi(cm),
 \qquad\text{number of raw covering rows}\le\pi(cm).
 \label{eq:tr-active-counts}
\end{equation}

For any nonnegative raw dual vector, all but at most $k(M)$ core weights
are at most $\eta=c_\delta/(k(M)+1)$, simply because the entire core has
at most $k(M)$ primes. Theorem~\ref{thm:cg-controlled} applies with
$k=k(M)$ and $R=R_B$. If $X$ is its raw load, the dual objective is
\[
 \sum_{a\in B}X(a)-\sum_{b\in I}(X(b)-1)_+.
\]
Subtracting the source hinge at threshold $2+\delta$ bounds this by
\[
 (2+\delta)m+B_{\delta,k(M)}(1+|R_B|).
\]
The bound is uniform over all dual vectors; strong LP duality and rounding
therefore give
\begin{equation}
 \tau(B,I)\le(2+\delta)m+B_{\delta,k(M)}
       +(B_{\delta,k(M)}+1)\pi(cm).
 \label{eq:tr-defect-round}
\end{equation}
There is no combination of independent covers in this step.

A failing subset again satisfies \eqref{eq:tr-minimal-large}; also
$m\le M-1$. Thus $\log m$ and $\log M$ are comparable for large $M$.
From \eqref{eq:tr-error-growth} and $k(M)=o(\log\log M)$,
\[
 B_{\delta,k(M)}\le(\log M)^{o(1)}=o(m).
\]
Chebyshev's bound now implies
\[
 (B_{\delta,k(M)}+1)\pi(cm)
   =O\!\left(\frac{(\log M)^{o(1)}}{\log m}\,m\right)
   =o(m).
\]
For all sufficiently large $M$, the errors in
\eqref{eq:tr-defect-round} sum to at most $\delta m$, contradicting the
choice of $B$. Since $\delta$ was arbitrary and $M\ge|A|+1$, the
uniform asymptotic conclusion follows.
\end{proof}

The corollary is an upper-bound theorem and does not assert a matching
lower coefficient. Its core-size restriction is much stronger than the
$O(\sqrt{n/\log n})$ core still permitted by the general necessary
structure. It therefore does not settle unrestricted Erd\H{o}s~708.

\section{An integer cover for a shared polylogarithmic core}
\label{sc:section}

The preceding fractional argument has an integer consequence for cores
of size \(o(\log\log M)\). A different construction permits a much
larger common core. It uses the source factorization itself: outside
prime powers are covered first, regular core parts are assigned to
disjoint target bins, and the remaining core demands are repaired
globally. No separation between prime sizes is assumed.

\begin{theorem}[A shared polylogarithmic core]
\label{thm:sc-polylog}
Fix \(0<a<2\). For every \(\varepsilon>0\), all sufficiently large
\(M\) have the following property. Suppose
\(A\subseteq\{2,\ldots,M\}\) contains \(n\) distinct integers and
admits a designated prime set \(E\) with
\[
 |E|\le(\log M)^a
\]
such that each source has at most one distinct prime factor outside
\(E\). Arbitrary powers of that outside prime are allowed.
Then, for every positive interval \(I\) of \(M\) consecutive
integers,
\begin{equation}
 \tau(A,I)\le(2+\varepsilon)n.
 \label{sc:main-bound}
\end{equation}
The threshold is uniform in the actual core and outside primes, the
source set, its cardinality, and the interval shift. In particular,
\(\tau(A,I)\le(2+o(1))n\) uniformly over this class as \(n\to\infty\).
\end{theorem}

The core may vary with the instance. It is shared by the entire
source set: a mixed source has the unique form \(d r^e\), where
\(d\) is \(E\)-smooth, \(r\notin E\) is prime, and \(e\ge1\).
The proof below is entirely integral. It uses neither a minimal
counterexample nor fractional rounding.

\subsection{One target per outside-prime occurrence group}

\begin{lemma}[Nested prime-power matching]
\label{sc:outside-matching}
For a distinct source set \(A\subseteq[2,M]\) and a prime \(r\),
the source exponents \(v_r(a)>0\) can be assigned to distinct
targets \(b\in I\), each satisfying \(v_r(b)\ge v_r(a)\).
Consequently, when each source has at most one distinct prime
outside \(E\), there is a target set \(B_{\rm out}\) of size at
most \(n_{\rm out}\) covering the full demand of every outside
prime. Here \(n_{\rm out}\) counts all mixed sources.
\end{lemma}

\begin{proof}
For every \(j\ge1\), distinctness gives
\begin{equation}
 \#\{a\in A:v_r(a)\ge j\}
 \le\left\lfloor\frac M{r^j}\right\rfloor
 \le\#\{b\in I:r^j\mid b\}.
 \label{sc:nested-hall}
\end{equation}
The allowed target sets are nested. Order the source exponents
decreasingly and assign them in that order. At a current exponent
\(e\), its allowed target set has at least as many positions as
there are source exponents at least \(e\). The previous assignments
have used fewer positions than that number, so an allowed unused
target exists. This proves the assertion for one prime.

Perform this matching separately for each outside prime \(r\), and
let \(B_r\) be the selected set. Take
\[
 B_{\rm out}=\bigcup_{r\notin E}B_r.
\]
For each \(r\), the union retains all of \(B_r\), and therefore
covers the entire \(r\)-demand. Different groups demand disjoint
prime coordinates; overlaps between their selected sets cause no
loss. Since each mixed source belongs to exactly one such group,
\(|B_{\rm out}|\le\sum_r|B_r|=n_{\rm out}\).
\end{proof}

\begin{proposition}[A finite certificate valid for every shift]
\label{sc:finite-certificate}
Under the one-distinct-outside-prime source condition, let \(d(a)\)
be the full \(E\)-smooth part of \(a\), and let
\(m_d=\#\{a\in A:d(a)=d\}\).
Choose any collection \(\mathcal D\) of occupied core parts \(d>1\).
Suppose the number of targets with exact core part \(d\) is at least
\(m_d\), for every shift and every \(d\in\mathcal D\). Then
\begin{equation}
 \tau(A,I)\le n_{\rm out}+\sum_{d\in\mathcal D}m_d
       +\sum_{\substack{a\in A\\d(a)>1,\ d(a)\notin\mathcal D}}
                    \Omega(d(a))
 \label{sc:finite-cover}
\end{equation}
for every shift. In particular, if the right side is at most \(2n\),
this certifies the exact \(2n\) bound for the given source set.
\end{proposition}

\begin{proof}
Use Lemma~\ref{sc:outside-matching} for all outside demands.
Choose \(m_d\) distinct targets from each selected exact-core-part
bin; these bins are disjoint. Their product covers all core demand
from those fibres. In the union of the two target sets, only the
remaining core units can be deficient. Repair them one at a time,
using unused targets guaranteed by full-interval supply. At most
the last sum in \eqref{sc:finite-cover} additions are needed.
Core part one requires no carrier.
\end{proof}

This certificate is finite and explicit. For \(T=M/d\),
inclusion--exclusion supplies the uniform bin lower bound
\[
 L_E(T)=
 \sum_{\substack{S\subseteq E\\|S|\ {\rm even}}}
      \left\lfloor\frac T{q_S}\right\rfloor
 -\sum_{\substack{S\subseteq E\\|S|\ {\rm odd}}}
      \left\lceil\frac T{q_S}\right\rceil,
 \qquad q_S=\prod_{p\in S}p,\quad q_\varnothing=1.
\]
Thus \(L_E(M/d)\ge m_d\) is sufficient. The simpler lower bound
\(T\prod_{p\in E}(1-1/p)-2^{|E|}\) is also valid.
The asymptotic proof below obtains better bin estimates when the
core grows; neither finite test is asserted to be necessary.

\subsection{Uniform estimates from the core cardinality}

For \(0<\sigma<1\), let
\[
 \Psi_E(x)=\#\{d\le x:d\text{ is }E\text{-smooth}\},
 \qquad
 {\cal Z}_E(\sigma)=\prod_{p\in E}(1-p^{-\sigma})^{-1}.
\]
One is included among the smooth integers. Rankin's elementary
inequality gives, for every \(x>0\),
\begin{equation}
 \Psi_E(x)\le x^\sigma{\cal Z}_E(\sigma).
 \label{sc:rankin}
\end{equation}
For \(x<1\) the left side is zero, so the same bound remains valid.

If \(E=\{p_1<\cdots<p_k\}\), then \(p_i\ge i+1\), and
\[
 \begin{split}
 \log{\cal Z}_E(\sigma)
 &\le\frac1{1-2^{-\sigma}}\sum_{p\in E}p^{-\sigma}\\
 &\le C_\sigma k^{1-\sigma}.
 \end{split}
\]
The empty Euler product is one. Fix, from now on,
\begin{equation}
 \max\{0,1-1/a\}<\sigma<\frac12.
 \label{sc:sigma-choice}
\end{equation}
Such a choice exists for every fixed \(0<a<2\). It yields
\begin{equation}
 {\cal Z}_E(\sigma)
 \le\exp\!\bigl(C_\sigma(\log M)^{a(1-\sigma)}\bigr)
 =M^{o(1)}
 \label{sc:euler-uniform}
\end{equation}
uniformly over all core prime values, since \(a(1-\sigma)<1\).

\begin{lemma}[Many targets with an exact core part]
\label{sc:bins}
Fix \(a>0\) and \(0<\theta<1\). There is \(c_a>0\) such that,
for all sufficiently large \(M\), every prime set
\(|E|\le(\log M)^a\) and every real interval \(J\) of length
\(M^\theta\le T\le M\) satisfy
\begin{equation}
 \#\{t\in J:(t,E^\Pi)=1\}
 \ge c_a\,\frac T{\log\log M}.
 \label{sc:bin-density}
\end{equation}
The large-\(M\) threshold may depend on \(a,\theta\); the estimate
is uniform in all prime values and shifts.
\end{lemma}

\begin{proof}
We use the inherited uniform interval sieve \cite{Synthesis},
derived from the fundamental lemma \cite{Lagarias,GranvilleSound}.
Precisely, if \(Q\) is a subset of primes at most \(Y\) and an
interval has real length \(T\ge Y^u\), then its count of integers
coprime to \(Q^\Pi\) differs from
\(T\prod_{q\in Q}(1-1/q)\) by at most \(e(u)\) times that
quantity, where \(e(u)\to0\) uniformly in \(Q\) and the shift.
The explicit error is stated in Proposition~\ref{cg:conditional-input}
and in the sieve input of the separated-support section.
We also use
\(\prod_{p\le Y}(1-1/p)\gg1/\log Y\)
\cite{GranvilleSound,TaoMertens}.

Set
\[
 Y=(\log M)^{2a+3},\qquad
 E_{\rm small}=E\cap[2,Y],\qquad E_{\rm large}=E\setminus E_{\rm small}.
\]
Since \(T\ge M^\theta\), the ratio \(\log T/\log Y\) tends to
infinity uniformly. For large \(M\), the small-prime sieve therefore
supplies at least
\[
 \frac12 T\prod_{q\in E_{\rm small}}(1-1/q)
 \gg\frac T{\log Y}
\]
integers coprime to \(E_{\rm small}\).
Requiring coprimality to the remaining core primes discards at most
\[
 \sum_{q\in E_{\rm large}}\left(\frac Tq+1\right)
 \le\frac{T|E|}{Y}+|E|
 =o(T/\log Y).
\]
Indeed
\[
 \frac{|E|\log Y}{Y}\to0,\qquad
 \frac{|E|\log Y}{T}\to0
\]
uniformly at the stated range of \(T\). Subtracting this upper
bound leaves a fixed positive multiple of \(T/\log Y\).
Finally \(\log Y=(2a+3)\log\log M\), proving the lemma.
\end{proof}

If \(d\) is \(E\)-smooth, the target positions with full
\(E\)-smooth part exactly \(d\) are
\[
 \{b\in I:d\mid b,\ (b/d,E^\Pi)=1\}.
\]
They are counted by applying Lemma~\ref{sc:bins} to the scaled
real interval \(I/d\), of length \(M/d\). Different \(d\)'s
give disjoint target bins.

\subsection{The three-stage construction}

\begin{proof}[Proof of Theorem~\ref{thm:sc-polylog}]
Keep the fixed \(\sigma\) from \eqref{sc:sigma-choice}, and put
\begin{equation}
 \theta=\frac{1/2-\sigma}{2(1-\sigma)},\qquad
 Z=M^\theta,\qquad
 \beta=\sigma+\theta(1-\sigma)=\frac14+\frac{\sigma}{2}<\frac12.
 \label{sc:cutoff}
\end{equation}
The parameters depend only on \(a\).
For every mixed source write \(a=d f\), where \(d\) is its full
\(E\)-smooth part and \(f=r^e\) is its outside prime-power part.
Call it regular when \(M/d\ge Z\), and bad otherwise.
All pure \(E\)-smooth sources are also called bad.
Let \(n_{\rm reg}\) be the number of regular sources, and let
\(n_{\rm out}\) count all mixed sources, including bad ones.

\emph{The bad sources have negligible total core demand.}
For a bad mixed source, \(d>M/Z\) and \(df\le M\), so \(f<Z\).
Enlarge the possible outside prime powers to all integers
\(2\le f<Z\). Equations \eqref{sc:rankin}--\eqref{sc:euler-uniform}
give
\begin{equation}
 \begin{split}
 |A_{\rm bad}|
 &\le\Psi_E(M)+\sum_{2\le f<Z}\Psi_E(M/f)\\
 &\le{\cal Z}_E(\sigma)M^\sigma
                    \left(1+\sum_{2\le f<Z}f^{-\sigma}\right)\\
 &\le C_\sigma{\cal Z}_E(\sigma)M^\sigma Z^{1-\sigma}
 =M^{\beta+o(1)}.
 \end{split}
 \label{sc:bad-count}
\end{equation}
Dropping the restriction \(d>M/Z\) only enlarged the count.
Write \(\Omega_E(a)=\sum_{p\in E}v_p(a)\).
The total bad core demand is at most
\begin{equation}
 \sum_{a\in A_{\rm bad}}\Omega_E(a)
 \le |A_{\rm bad}|\log_2M
 =M^{\beta+o(1)}
 =o(\sqrt M).
 \label{sc:bad-demand}
\end{equation}
All errors are uniform in the admitted source data.

\emph{First select targets for every outside-prime demand.}
Lemma~\ref{sc:outside-matching} gives a set \(B_{\rm out}\)
of at most \(n_{\rm out}\) positions. Its product covers the
entire demand of each prime outside \(E\), including those in
bad mixed sources.

\emph{Next select simultaneous carriers for the regular core parts.}
Group regular sources by their exact core part \(d\). Put
\(T_d=M/d\), and let \(m_d\) be the fibre size.
Distinct sources in the same fibre have distinct outside prime
powers \(f\le T_d\). If \({\cal P}(T)\) counts prime powers
\(p^j\le T\), \(j\ge1\), then
\[
 m_d\le{\cal P}(T_d),\qquad
 {\cal P}(T)\le\pi(T)+
       \sum_{2\le j\le\log_2T}\pi(T^{1/j})
 =O(T/\log T).
\]
Here Chebyshev bounds the first term; the remaining terms total
at most \(\sqrt T\log_2T=o(T/\log T)\).
On every occupied regular fibre \(T_d\ge M^\theta\), so
Lemma~\ref{sc:bins} gives an exact-core-part bin of size
\[
 N_d\ge c_a\,\frac{T_d}{\log\log M}.
\]
Consequently
\[
 \frac{m_d}{N_d}
 =O_a\!\left(\frac{\log\log M}{\log T_d}\right)
 =O_a\!\left(\frac{\log\log M}{\log M}\right)\longrightarrow0.
\]
Thus \(N_d\ge m_d\) simultaneously in all fibres for large
\(M\). Choose \(m_d\) distinct targets from each bin. Since the
bins are disjoint, this yields a set \(B_{\rm core}\) of exactly
\(n_{\rm reg}\) positions with
\begin{equation}
 \prod_{a\in A_{\rm reg}}d(a)\ \bigm|\ \prod_{b\in B_{\rm core}}b.
 \label{sc:core-carriers}
\end{equation}
This is one joint covering of all regular core multiplicities.

\emph{Finally repair the remaining core units globally.}
Take \(B_0=B_{\rm out}\cup B_{\rm core}\). Its cardinality is
at most \(n_{\rm out}+n_{\rm reg}\).
For every outside prime it retains all of \(B_{\rm out}\);
for every core prime it retains all of \(B_{\rm core}\).
Therefore the only possible deficits are core units from the
bad sources, whose total is bounded in \eqref{sc:bad-demand}.
If a prime remains deficient, full-interval supply
\(\prod A\mid M!\mid\prod I\) guarantees an unused target
divisible by that prime. Add such a target. Each addition reduces
the total positive deficit by at least one and never increases
another deficit. Hence
\begin{equation}
 \begin{split}
 \tau(A,I)
 &\le n_{\rm out}+n_{\rm reg}
                 +\sum_{a\in A_{\rm bad}}\Omega_E(a)\\
 &\le2n+M^{\beta+o(1)}.
 \end{split}
 \label{sc:constructive-bound}
\end{equation}
Overlaps between the first two selected sets are harmless, since
their demanded prime families are disjoint. The last repair uses
the accumulated cover, rather than an independent cover of the
bad sources.

To finish uniformly over all cardinalities, if \(n\le\sqrt M/4\),
the ambient quadratic theorem \cite[Appendix A]{Wang} gives
\(\tau(A,I)\le2n\). Otherwise \(n>\sqrt M/4\), and
\[
 \frac{M^{\beta+o(1)}}n
 \le4M^{\beta-1/2+o(1)}\longrightarrow0.
\]
This proves \eqref{sc:main-bound} for all sufficiently large
\(M\), uniformly over the claimed class and shifts.
Finally \(n\to\infty\) forces \(M\ge n+1\to\infty\),
giving the stated uniform asymptotic conclusion.
\end{proof}

\subsection{An explicit core bound approaching the squared logarithm}

\begin{corollary}[A quantified near-square core]
\label{cor:sc-near-square}
For all sufficiently large \(M\), suppose
\begin{equation}
 |E|\le10^{-6}
       \left(\frac{\log M}{\log\log M}\right)^2
 \label{sc:near-square-size}
\end{equation}
and each distinct source has at most one distinct prime outside
\(E\), with arbitrary powers allowed. Then
\(\tau(A,I)\le(2+o(1))|A|\), uniformly in the source set,
cardinality, prime values, and shift. More precisely, with
\(L=\log M\) and \(l=\log L\), the direct construction gives
\begin{equation}
 \tau(A,I)\le2|A|+
 O\!\left(\sqrt M\,L\exp\left(-\frac{2L}{5l}\right)\right).
 \label{sc:near-square-error}
\end{equation}
The absolute numerical constant in \eqref{sc:near-square-size}
is a sufficient choice and is not asserted to be optimal.
\end{corollary}

\begin{proof}
We check the varying parameters in the same construction explicitly,
rather than substituting them into a theorem for fixed parameters.
Assume \(l\ge4\), and set
\[
 \delta=\frac1l,\qquad
 \sigma=\frac12-\delta\in[1/4,1/2),\qquad
 \theta=\frac{\delta}{2(1-\sigma)}=\frac1{l+2},\qquad
 \beta=\frac12-\frac1{2l}.
\]
Let \(k=|E|\). The Euler estimate is uniform in this range of
\(\sigma\):
\[
 \begin{split}
 \log{\cal Z}_E(\sigma)
 &\le\frac1{1-2^{-\sigma}}
                       \sum_{i=1}^k(i+1)^{-\sigma}\\
 &\le\frac{2}{1-2^{-1/4}}\,k^{1-\sigma}
 <13k^{1-\sigma}.
 \end{split}
\]
The empty core is immediate. For \(k\ge1\), the bound \(k\le L^2\)
implies \(k^\delta\le e^2\). Therefore
\begin{equation}
 \log{\cal Z}_E(\sigma)
 <13e^2\sqrt k<100\sqrt k\le\frac{L}{10l}.
 \label{sc:near-square-euler}
\end{equation}
Here \(\sqrt k\le10^{-3}L/l\) follows from
\eqref{sc:near-square-size}.

Define the regular mixed sources by \(M/d\ge M^\theta\), as
before. The sum in \eqref{sc:bad-count} is bounded uniformly by
\(3(M^\theta)^{1-\sigma}\), since \(1/(1-\sigma)\le2\).
Combining this with \eqref{sc:near-square-euler} gives
\[
 |A_{\rm bad}|
 \le3\sqrt M\exp\left(-\frac{L}{2l}+\frac{L}{10l}\right)
 =3\sqrt M\exp\left(-\frac{2L}{5l}\right).
\]
Consequently the total bad core demand is at most this bound times
\(\log_2M\), which is \(o(\sqrt M)\).

It remains to check the target bins with the shrinking \(\theta\).
Take \(Y=L^7\). Every occupied regular fibre has
\[
 T\ge M^\theta=\exp\left(\frac L{l+2}\right),\qquad
 \frac{\log T}{\log Y}\ge\frac{L}{7l(l+2)}\longrightarrow\infty.
\]
The uniform small-prime sieve still gives a fixed positive multiple
of \(T/\log Y\) positions coprime to \(E\cap[2,Y]\).
Deleting multiples of the remaining core primes costs at most
\[
 \frac{Tk}{Y}+k\le\frac T{L^5}+L^2=o(T/\log Y).
\]
Thus exact-core-part bins have size \(\gg T/l\).
The number of possible prime-power remainders remains
\(O(T/\log T)\), and its ratio to the bin size is
\[
 O\left(\frac l{\log T}\right)
 =O\left(\frac{l(l+2)}L\right)\longrightarrow0.
\]
All bins can therefore be assigned as before. The outside-prime
matching and the final global repair do not depend on
\(\sigma,\theta\). They yield \eqref{sc:near-square-error}.

If \(|A|\le\sqrt M/4\), use the exact quadratic theorem.
Otherwise divide the error in \eqref{sc:near-square-error} by
\(|A|>\sqrt M/4\). The relative error is
\(O(L\exp(-2L/(5l)))=o(1)\), uniformly, proving the corollary.
\end{proof}

For every fixed \(a<2\), one has
\((\log M)^a\le10^{-6}(\log M/\log\log M)^2\) eventually.
Thus the quantified corollary includes the fixed-exponent theorem,
while its proof checks the additional uniformity needed when the
effective exponent approaches two.

\begin{corollary}[A stronger core order from Chebyshev's estimate]
\label{cor:sc-chebyshev-core}
There is an absolute constant \(c_0>0\) such that the same uniform
integer conclusion holds under the source factorization hypothesis
of Theorem~\ref{thm:sc-polylog}, whenever
\begin{equation}
 |E|\le c_0\,\frac{(\log M)^2}{\log\log M}.
 \label{sc:chebyshev-core-size}
\end{equation}
For sufficiently large \(M\), the construction also satisfies the
error bound \eqref{sc:near-square-error}.
The existence of a positive absolute \(c_0\) is asserted; no optimal
constant is claimed.
\end{corollary}

\begin{proof}
Chebyshev's estimate implies \(p_i\gg i\log(i+1)\) for the
\(i\)-th prime \cite{GranvilleSound}. Uniformly for
\(1/4\le\sigma\le1/2\), this gives
\begin{equation}
 \log{\cal Z}_E(\sigma)
 \le C\,\frac{k^{1-\sigma}}{(\log(k+1))^\sigma},
 \qquad k=|E|\ge1,
 \label{sc:chebyshev-euler}
\end{equation}
with an absolute \(C\). To check uniformity, split the sum of
\(p_i^{-\sigma}\) at \(i=\sqrt k\). The first part is
\(O(k^{(1-\sigma)/2})\). In the second part
\(\log(i+1)\gg\log(k+1)\), giving
\(O(k^{1-\sigma}/(\log(k+1))^\sigma)\).
The first part is absorbed in the second bound uniformly in
this compact range of \(\sigma\); its relative ratio is at most
a constant times \(\sqrt{\log(k+1)}/k^{1/4}\).
The Euler logarithm introduces only the bounded factor
\((1-2^{-\sigma})^{-1}\). Enlarging \(C\) handles small \(k\).

Use \(L=\log M\), \(l=\log L\), and the same
\(\sigma=1/2-1/l\) as in Corollary~\ref{cor:sc-near-square}.
Put \(K=\lfloor c_0L^2/l\rfloor\).
The function \(k^{1-\sigma}/(\log(k+1))^\sigma\) is increasing:
its logarithmic derivative multiplied by \(k\) is
\[
 (1-\sigma)-
 \frac{\sigma k}{(k+1)\log(k+1)}
 \ge1-2\sigma\ge0.
\]
For every fixed \(c_0>0\), as \(M\to\infty\),
\[
 \begin{split}
 \frac{K^{1-\sigma}}{(\log(K+1))^\sigma}
 &=\sqrt{\frac K{\log(K+1)}}\,
   \exp\left(\frac{\log K+\log\log(K+1)}l\right)\\
 &=\left(e^2\sqrt{\frac{c_0}{2}}+o(1)\right)\frac L l.
 \end{split}
\]
Choose \(c_0\le1\) so small that
\(Ce^2\sqrt{c_0/2}<1/20\). Then
\eqref{sc:chebyshev-euler}, for every \(k\le K\), is at most
\(L/(10l)\) once \(M\) is large.
The empty core causes no error.

This is precisely the Euler bound
\eqref{sc:near-square-euler} needed in the preceding proof.
Also \(k\le L^2\) for large \(M\), so its \(Y=L^7\)
sieve calculation, large-prime deletion bound, prime-power fibre
comparison, and bad-demand estimate remain valid without change.
The same construction and the quadratic gate finish the proof.
\end{proof}

\subsection{Scope of the larger shared-core class}

For example, a designated core of size at most \((\log M)^{1.99}\)
is permitted, with any number of possible outside primes and no
prime-size gap. The source restriction still concerns the number
of \emph{distinct} outside primes, so their arbitrary powers remain
fully included.

The two conditions on the Rankin exponent are
\[
 a(1-\sigma)<1,\qquad \sigma<1/2.
\]
Their simultaneous use explains the endpoint \(a<2\) in the
fixed-parameter argument. Corollaries~\ref{cor:sc-near-square}
and \ref{cor:sc-chebyshev-core} check specified approaches to
that endpoint. These results do not
establish a core of size \((\log M)^2\), an arbitrary
instance-dependent exponent tending to two, or an optimal endpoint.

The unrestricted minimal-obstruction core can still have size
\(O(\sqrt{n/\log n})\), much larger than a fixed power of
\(\log M\) in the subquadratic range. Thus
Theorem~\ref{thm:sc-polylog} enlarges a concrete source class
without resolving the unrestricted asymptotic or exact conjecture.
Its integer construction is independent of the controlled-weight
hinge bound: the latter remains useful for restricting dual vectors
inside a larger designated core.

\section{Disjoint support components: a dependent consequence and scope check}
\label{bs:section}

The aggregate-core estimates suggest covering prime-disjoint source
components separately. That argument is valid, but for components with
$o(\log M)$ primes it does not produce a previously untreated integer
class: an elementary smooth-number count and the inherited quadratic
theorem already give an eventual exact $2n$ bound. We record the stronger
dependent consequence to make this limitation explicit.

\begin{lemma}[A uniform count on a small prime support]
\label{bs:rankin}
For every finite prime set $Q$ with $|Q|\le k$, the number $\Psi_Q(M)$
of $Q$-smooth positive integers at most $M$ satisfies
\begin{equation}
 \Psi_Q(M)\le M^{1/4}\exp(Ck^{3/4}),
 \qquad C=\frac4{3(1-2^{-1/4})}.
 \label{bs:smooth-count}
\end{equation}
The bound is uniform in the actual prime values.
\end{lemma}

\begin{proof}
Rankin's inequality at exponent $1/4$ gives
\[
 \Psi_Q(M)\le
 M^{1/4}\sum_{\substack{h\ge1\\h\ Q\text{-smooth}}}h^{-1/4}
 =M^{1/4}\prod_{p\in Q}(1-p^{-1/4})^{-1}.
\]
If the primes are listed as $p_1<\cdots<p_f$, then $p_j\ge j+1$.
Using $-\log(1-x)\le x/(1-x)$, we obtain
\[
 \begin{split}
 \log\prod_{p\in Q}(1-p^{-1/4})^{-1}
 &\le\frac1{1-2^{-1/4}}\sum_{j=1}^f(j+1)^{-1/4}\\
 &\le\frac4{3(1-2^{-1/4})}k^{3/4}.
 \end{split}
\]
The empty support also satisfies the bound, since it contains only
the smooth integer one.
\end{proof}

\begin{corollary}[An inherited exact bound for disjoint small supports]
\label{bs:exact}
Let $k(M)=o(\log M)$. Suppose the prime support of a distinct source set
$A\subseteq[2,M]$ can be partitioned into disjoint blocks
$P_1,\ldots,P_s$, each of size at most $k(M)$, and every source has
\emph{all} its prime factors in a single block. Then, for all sufficiently
large $M$, uniformly in the partition, the prime values, and every
positive $M$-term interval,
\[
 \tau(A,I)\le2|A|.
 \]
\end{corollary}

\begin{proof}
Group the sources as $A_i$, supported on $P_i$, and discard empty
groups. Let $n_i=|A_i|$. Lemma~\ref{bs:rankin} gives, uniformly in $i$,
\[
 n_i\le M^{1/4}\exp(Ck(M)^{3/4})
      =M^{1/4+o(1)}<\sqrt M/4
\]
for large $M$. Hence $M\ge16n_i^2$ for every group.
The ambient quadratic theorem of \cite{Wang} supplies a cover $B_i$
of $A_i$ in the original interval $I$ with $|B_i|\le2n_i$.

The union $B=\bigcup_iB_i$ covers all of $A$.
Indeed, each demanded prime belongs to a unique block. Its entire
source demand is supplied by that block's cover, which is contained
in $B$. Possible overlaps of different $B_i$ therefore cause no loss:
there is no additional source demand at that prime from another
block. Finally,
$|B|\le\sum_i|B_i|\le2\sum_i n_i=2|A|$.
\end{proof}

For comparison, the aggregate-core bound with no outside primes gives
\[
 \mathcal L_{\rm raw}(A_i,I)
 \le(2+\delta)n_i+B_{\delta,k},\qquad
 B_{\delta,k}\le D_\delta^k,
\]
and rounding adds at most $k$ positions. Applying this only to groups
with $n_i>\sqrt M/4$, while using the exact quadratic bound for smaller
groups, would give the valid estimate
\[
 \tau(A,I)\le(2+\delta)|A|
       +\frac{4|A|}{\sqrt M}(B_{\delta,k}+k).
\]
For $k=o(\log M)$ its last term is $o(|A|)$.
However, the smooth-number argument shows that eventually every group
is already small, so the fractional route is unnecessary for this class.

The component hypothesis is stronger than a bound on the number of
distinct primes in each individual source. Different sources can connect
many prime values through overlaps. It also differs from a shared-core
geometry, where common primes occur across many source groups; in that
case the prime-disjoint union argument is unavailable.
The auxiliary weighted block-gluing theorem permits such a shared load
only when its weights are diffuse, and it does not remove the arbitrary
shared-core obstruction.

\section{Separated prime supports and discounted certificates}
\label{ss:section}

This section proves a uniform coefficient-two asymptotic bound for sources
whose prime factors lie in two separated ranges. The large-prime range may
start at $M^\alpha$ for any fixed $\alpha>0$. A quantitative extension also
allows a specified exponent tending to zero. The intermediate prime range
remains excluded.

All logarithms are natural unless indicated otherwise. For a prime set $Q$,
write $Q^\Pi=\prod_{q\in Q}q$ and
$\vartheta(Q)=\prod_{q\in Q}(1-1/q)$, with empty products equal to one.
An integer is $Q$-smooth when all its prime factors belong to $Q$;
one is included. Intervals at real lengths are half-open intervals
$(x,x+T]$ with $x\ge0$, with their integer points understood.
The ambient integer $M$ is retained when sources are deleted.

\subsection{Precise analytic inputs}

We recall the following previously proved statements, including the
coefficient in the diffuse comparison. Their proofs, with partial
prime-power carriers and uniform interval errors, are in \cite{Synthesis}.
The new arguments below use the coefficients as stated.

\begin{lemma}[Uniform sieve and diffuse comparison inputs]
\label{ss:inputs}
There is an absolute $C>0$ such that
\[
 e(u)=C\bigl((u/2)^{-u/4}+2^{1-u/2}\bigr),\qquad u\ge4,
\]
is decreasing to zero and has the following property. If $Q$ is any
subset of the primes at most $z\ge2$ and an interval $J$ has real length
$T\ge z^u$, then
\[
 \left|\#\{m\in J:(m,Q^\Pi)=1\}-T\vartheta(Q)\right|
 \le e(u)T\vartheta(Q).
 \tag{S}
\]
For weights $0\le w_p\le\eta$, put
$V(m)=\sum_pw_pv_p(m)$ and $\mu=\sum_pw_p/(p-1)$.
Let $\ell>\eta$, $r=\lfloor\ell/\eta\rfloor\ge8$,
$j=\lfloor r/2\rfloor$, and $\kappa=r-j$.
Whenever $e(r)<1$, every level $s\ge0$ and every pair of a prefix and
an interval of the same real length $T$ satisfy
\begin{equation}
 \#\{a\le T:V(a)>s+\ell\}
 \le
 \frac{(1+e(\kappa))\mu+\eta2^{1-j}}
      {(\ell-\eta)(1-e(r))}
 \#\{b\in J:V(b)>s\}.
 \label{ss:diffuse-ratio}
\end{equation}
Finally, if the support is contained in $p\le M$, then
\begin{equation}
 \sum_{a=1}^M V(a)\ge M(\mu-3\eta).
 \label{ss:small-mean-error}
\end{equation}
\end{lemma}

For clarity, the carrier proof of \eqref{ss:diffuse-ratio} bounds the
source count by a common carrier sum times
$((1+e(\kappa))\mu+\eta2^{1-j})/(\ell-\eta)$, and bounds the target
count below by that sum times $1-e(r)$. Thus retaining their quotient
gives exactly the displayed statement.
We also use the classical uniform bounds
\begin{equation}
 \pi(T)\le C_0T/\log T,\qquad
 \vartheta(Q)\ge c_0/\log Y
 \quad(Q\subseteq\{p:p\le Y\},\ Y\ge2),
 \label{ss:prime-inputs}
\end{equation}
with absolute positive constants. These follow from Chebyshev and
Mertens estimates \cite{GranvilleSound,TaoMertens}.

\subsection{A strict budget for arbitrary pure-core sources}

\begin{proposition}[Uniform pure-core contraction]
\label{ss:pure-core}
Let $Y=Y(M)\ge2$ be prescribed with $\log Y=o(\sqrt{\log M})$.
For every $\varepsilon>0$ there is $\rho_\varepsilon<1$ such that,
for all sufficiently large $M$, the following holds uniformly in
$Q\subseteq\{p:p\le Y\}$, weights $0\le w_q\le1$, distinct
$Q$-smooth sources $A_Q\subseteq[2,M]$, and positive $M$-term intervals:
\begin{equation}
 \sum_{a\in A_Q}(V(a)-2-\varepsilon)_+
 \le \rho_\varepsilon
       \sum_{b\in I}(V(b)-1)_+,
 \qquad V(m)=\sum_{q\in Q}w_qv_q(m).
 \label{ss:pure-contraction}
\end{equation}
No bound on the number of prime factors of a pure-core source is assumed.
\end{proposition}

\begin{proof}
It suffices to treat $0<\varepsilon\le1$.
Set $\ell=1+\varepsilon$ and
\[
 \rho_\varepsilon=\frac{1+\varepsilon/2}{1+\varepsilon}<1.
\]
Choose $\eta>0$ so small that $\eta\le\varepsilon/24$, $r\ge8$,
$e(r)<1$, and
\begin{equation}
 \frac{(1+e(\kappa))(1+\varepsilon/4)+\eta2^{1-j}}
      {(\ell-\eta)(1-e(r))}
 \le\rho_\varepsilon,
 \label{ss:strict-choice}
\end{equation}
where $r,j,\kappa$ are as in Lemma~\ref{ss:inputs}.
Such a choice exists: as $\eta\downarrow0$, the left side tends to
$(1+\varepsilon/4)/(1+\varepsilon)$, strictly below $\rho_\varepsilon$.
All these parameters depend only on $\varepsilon$.

We first record a uniform smooth-number bound. Let $\Psi(M,Y)$ count
positive integers at most $M$ all of whose prime factors are at most $Y$.
Write $N=\log M$, $y=\log Y$, and $\sigma=1/(4y)$.
For $p\le Y$, put $z_p=p^{-1+\sigma}$. Then
$z_p\le e^{1/4}/p$ and $z_p/(1-z_p)\le4/p$. Consequently
\[
 \prod_{p\le Y}(1-p^{-1+\sigma})^{-1}
 \le\exp\left(4\sum_{p\le Y}\frac1p\right)\le Y^4.
\]
Rankin's inequality therefore gives
\begin{equation}
 \Psi(M,Y)\le
 M\exp\left(-\frac{N}{4y}+4y\right).
 \label{ss:smooth-count}
\end{equation}
For every fixed $B\ge0$, it follows that
\begin{equation}
 \frac{\Psi(M,Y)\log_2M}{M/Y^B}
 \le(\log_2M)
       \exp\left(-\frac N{4y}+(B+4)y\right)\longrightarrow0.
 \label{ss:smooth-small}
\end{equation}
Indeed $y=o(\sqrt N)$ makes $N/y$ dominate both $y$ and $\log N$.
The bounds are independent of the prime subset and the weights.

Put $\mu=\sum_{q\in Q}w_q/(q-1)$ and
$H(V)=\sum_{b\in I}(V(b)-1)_+$.
There are three cases.

\emph{All weights are at most $\eta$, and $\mu\le1+\varepsilon/4$.}
Equations \eqref{ss:diffuse-ratio} and \eqref{ss:strict-choice} give
a level-count ratio at most $\rho_\varepsilon$.
Integrating over $s\ge1$ yields
\[
 \sum_{a\le M}(V(a)-2-\varepsilon)_+
 \le\rho_\varepsilon H(V).
\]
Distinctness permits restriction to $A_Q$.
This case includes the empty support and zero weights.

\emph{All weights are at most $\eta$, and $\mu\ge1+\varepsilon/4$.}
For large $M$, all $q\le Y$ are below $M$.
Every interval prime-power count is at least its prefix count, so
\[
 \begin{split}
 H(V)&\ge\sum_{b\in I}V(b)-M
 \ge M(\mu-3\eta-1)
 \ge\varepsilon M/8.
 \end{split}
\]
For $a\le M$, the source load satisfies
$V(a)\le\Omega(a)\le\log_2M$. The source hinge is at most
$\Psi(M,Y)\log_2M=o(M)$ by
\eqref{ss:smooth-small} with $B=0$. Its ratio to $H(V)$ is eventually
at most $\rho_\varepsilon$.

\emph{Some weight $w_q$ is at least $\eta$.}
Put $j_q=\lceil2/w_q\rceil$ and $B=\lceil2/\eta\rceil$.
Every multiple of $q^{j_q}$ has $V\ge2$, while
$q^{j_q}\le Y^B$ and $M/Y^B\to\infty$. Hence
\[
 H(V)\ge\left\lfloor\frac M{q^{j_q}}\right\rfloor
 \ge\frac M{2Y^B}
\]
for large $M$. The source hinge is again at most
$\Psi(M,Y)\log_2M$, which is $o(M/Y^B)$ by
\eqref{ss:smooth-small}. Its ratio to $H(V)$ is eventually at most
$\rho_\varepsilon$.
The thresholds in all three cases are uniform, proving
\eqref{ss:pure-contraction}.
\end{proof}

\subsection{A pointwise certificate with controlled multiplicity}

Let $Q,R_0$ be disjoint finite prime sets, and let $\mathcal D$ be a finite
set of integers $d\le M$ with the following properties:
\[
 d\text{ is }(Q\cup R_0)\text{-smooth},\qquad
 \Omega_{R_0}(d)\le t,\qquad
 r\mid d,\ r\in R_0\ \Longrightarrow\ r\le T_d:=M/d.
 \tag{D}
\]
Here $t\ge0$ is an integer, and $\Omega_{R_0}$ counts occurrences with
multiplicity. The number of $Q$-prime occurrences is unrestricted.
For an integer $L\ge1$, set $K_d(b)=0$ when $d\nmid b$, and otherwise set
\begin{equation}
 K_d(b)=\mathbf1_{(b/d,Q^\Pi)=1}
       \left(1-\frac1L
          \sum_{\substack{r\in R_0\\r\le T_d}}\mathbf1_{r\mid b/d}\right).
 \label{ss:discounted}
\end{equation}
Define
\begin{align}
 C_1(t,L)&=\binom{2t+L-1}{t},\notag\\
 C_2(t,L)&=(t+L)\sum_{j=0}^{L-1}
                  \binom{t+L-1}{j}(t+1)^j,\notag\\
 C_*(t,L)&=\min\{C_1(t,L),C_2(t,L)\}.
 \label{ss:multiplicity}
\end{align}

\begin{lemma}[Discounted pointwise capacity]
\label{ss:capacity}
Let $X(m)=\sum_pw_pv_p(m)$ have nonnegative weights on any finite prime
support containing $Q\cup R_0$, and put $h_d=(X(d)-1)_+$.
Then, for every positive integer $b$,
\begin{equation}
 \sum_{d\in\mathcal D}h_dK_d(b)
 \le C_*(t,L)(X(b)-1)_+.
 \label{ss:pointwise}
\end{equation}
\end{lemma}

\begin{proof}
Suppose a positive $K_d(b)$ exists, and choose the numerically least
positive carrier $d_0$. Then $T_0=M/d_0$ is the largest threshold
among positive carriers. Put
\[
 S=\{r\in R_0:r\le T_0,\ r\mid b\},\qquad k=|S|.
\]
Every member of $S$ divides $d_0$ or $b/d_0$.
The first class has at most $t$ distinct members.
The second has at most $L-1$ members, since positivity in
\eqref{ss:discounted} means that its integer sum is less than $L$.
Consequently $k\le t+L-1$.

Every positive carrier has the same $Q$-part, namely the full $Q$-smooth
part of $b$. By (D), its $R_0$-part is a monomial on $S$ of total
degree at most $t$. There are $\binom{k+t}{t}\le C_1(t,L)$ possible
monomials, including the constant monomial. For $k=0$ this count is one.
The common $Q$-part and the monomial determine the carrier.

For the sharper count, order $S$ increasingly.
For a positive carrier $d$, the primes in $S$ at most $T_d$ form a
prefix, with at most $k+1$ choices. Outside that prefix, their exponents
in $d$ are zero by (D). Inside it, all but at most $L-1$ primes must
occur to their full exponent $v_r(b)$: every deficiency is counted
by $r\mid b/d$ in \eqref{ss:discounted}.
Choose the exceptional primes and their exponents, each in
$\{0,\ldots,t\}$. All other exponents are fixed by $b$. Thus the number
of positive carriers is at most
\[
 (k+1)\sum_{j=0}^{L-1}\binom{k}{j}(t+1)^j\le C_2(t,L).
 \]
Choices violating another restriction only enlarge this upper bound.

For each positive term, $K_d(b)\le1$ and $d\mid b$ implies
$h_d\le(X(b)-1)_+$. All other weighted terms are nonpositive.
Either multiplicity bound therefore proves \eqref{ss:pointwise};
if no term is positive the inequality is immediate.
\end{proof}

\subsection{Uniform interval averages of the certificate}

\begin{lemma}[A positive average despite a large reciprocal sum]
\label{ss:average}
Suppose $Q\subseteq\{p:p\le Y\}$, all occupied carriers have $T_d\ge Z$,
and $Z\to\infty$. Allow $t,L,Y,Z$ to depend on $M$.
Assume uniformly over the admitted data that
\begin{equation}
 \frac{\mu_0}{L}\le\frac14,\qquad
 \mu_0=\sum_{r\in R_0}\frac1r,\qquad
 \log Y=o(\log Z),\qquad
 \frac{(\log Y)^2}{L\log Z}\longrightarrow0.
 \label{ss:average-conditions}
\end{equation}
Then, for all sufficiently large $M$, uniformly in the interval shift
and $d\in\mathcal D$,
\begin{equation}
 G_d:=\sum_{b\in I}K_d(b)
 \ge\frac12\vartheta(Q)T_d.
 \label{ss:positive-average}
\end{equation}
\end{lemma}

\begin{proof}
Let $C_Q(J)$ count integers in a real interval $J$ coprime to $Q^\Pi$.
For $T=M/d$ and $J=(x/d,(x+M)/d]$, exact division gives
\begin{equation}
 G_d=C_Q(J)-\frac1L
      \sum_{\substack{r\in R_0\\r\le T}}C_Q(J/r).
 \label{ss:average-identity}
\end{equation}
Division by $r\in R_0$ preserves the $Q$-coprimality condition.

Fix an absolute $\kappa\ge4$ with $e(\kappa)\le1/16$ and put
$U=T/Y^\kappa$. By \eqref{ss:average-conditions}, for large $M$,
$T\ge Y^\kappa$ and $U\ge\sqrt T$, uniformly for $T\ge Z$.
Apply (S) to the base count and every subtracted term with $r\le U$.
Their combined contribution is at least
\[
 \begin{split}
 \vartheta(Q)T\left[
  1-\frac{\mu_0}{L}
       -e(\kappa)\left(1+\frac{\mu_0}{L}\right)\right]
 \ge\frac{43}{64}\vartheta(Q)T.
 \end{split}
\]
For the remaining terms, discard coprimality and use
$C_Q(J/r)\le T/r+1$. Chebyshev's bound and partial summation give
\[
 \begin{split}
 \sum_{U<r\le T}\frac1r
 &\le\frac C{\log T}
       +C\int_U^T\frac{du}{u\log u}\\
 &\le C'\frac{1+\kappa\log Y}{\log T},
 \end{split}
\]
where the sums are enlarged to all primes. In the last inequality,
$U\ge\sqrt T$ and $\log(T/U)=\kappa\log Y$ were used.
Thus the remaining negative contribution in
\eqref{ss:average-identity} is
\[
 O\left(\frac{T(1+\kappa\log Y)}{L\log T}\right).
 \]
After division by $\vartheta(Q)T$, equation
\eqref{ss:prime-inputs} bounds this by
\[
 O_\kappa\left(\frac{(\log Y)^2}{L\log Z}\right)=o(1).
 \]
Here $\log Y\ge\log2$ absorbs the constant term in the numerator.
Subtracting this error from $43/64$ leaves at least $1/2$ for large
$M$, proving the lemma.
\end{proof}

\subsection{The general separated-support theorem}

\begin{theorem}[A uniform shared-hinge bound]
\label{ss:general}
Let $Y(M)\ge2$, $Z(M)\to\infty$, and integer functions $t(M)\ge0$,
$L(M)\ge1$ be prescribed. Assume
\begin{equation}
 \log Y=o(\sqrt{\log M}),\qquad
 \log Y=o(\log Z),\qquad
 \frac{(\log Y)^2}{L\log Z}\to0,\qquad
 C_*(t,L)\frac{\log Y}{\log Z}\to0.
 \label{ss:general-conditions}
\end{equation}
Let $Q\subseteq\{p:p\le Y\}$ and $R\subseteq\{p:p\ge Z\}$ be finite
prime sets. Let $A\subseteq[2,M]$ be distinct sources supported on
$Q\cup R$. For each source containing an $R$-prime, put
\[
 p(a)=P^+(a),\qquad d(a)=a/p(a).
\]
Let $\mathcal D$ be the distinct quotients and $R_0$ the union of their
$R$-prime supports. Suppose
\[
 \Omega_R(d)\le t\quad(d\in\mathcal D),
 \qquad \sum_{r\in R_0}\frac1r\le L/4.
\]
For every $\varepsilon>0$, all sufficiently large $M$ then satisfy,
uniformly over these data, all shifts, and all weights in $[0,1]$,
with $X(m)=\sum_{p\in Q\cup R}w_pv_p(m)$,
\begin{equation}
 \sum_{a\in A}(X(a)-2-\varepsilon)_+
 \le \sum_{b\in I}(X(b)-1)_+.
 \label{ss:shared-hinge}
\end{equation}
Consequently $\mathcal L_{\rm raw}(A,I)\le(2+\varepsilon)|A|$.
\end{theorem}

\begin{proof}
For large $M$, $Y<Z$, so $Q$ and $R$ are disjoint and the largest prime
of a mixed source belongs to $R$. Write $A_Q$ for the pure $Q$-sources
and $A_R=A\setminus A_Q$. If $A_R$ is empty, use
Proposition~\ref{ss:pure-core}.

Otherwise every $d=d(a)$ has $T_d\ge p(a)\ge Z$.
Every $R$-prime dividing $d$ is at most $p(a)\le T_d$, so (D) holds.
Let $m_d$ be the number of source preimages of $d$. Distinct preimages
have the form $dp(a)$ for distinct primes $p(a)\le T_d$, hence
$m_d\le\pi(T_d)$. Lemma~\ref{ss:average} and
\eqref{ss:prime-inputs} give
\begin{equation}
 \frac{m_d}{G_d}
 \le C_3\frac{\log Y}{\log Z},
 \label{ss:fibre-ratio}
\end{equation}
with an absolute constant, uniformly over all occupied fibres.

Since $w_{p(a)}\le1$,
$(X(a)-2-\varepsilon)_+\le(X(d)-1)_+=h_d$.
Using positivity of $G_d$, then Lemma~\ref{ss:capacity}, gives
\[
 \begin{split}
 \sum_{a\in A_R}(X(a)-2-\varepsilon)_+
 &\le\sum_d m_dh_d\\
 &\le C_3\frac{\log Y}{\log Z}\sum_dG_dh_d\\
 &\le C_3C_*(t,L)\frac{\log Y}{\log Z}
                  \sum_{b\in I}(X(b)-1)_+.
 \end{split}
\]
The last coefficient tends to zero.
For $V$ consisting only of the $Q$ weights,
Proposition~\ref{ss:pure-core} gives a uniform $\rho_\varepsilon<1$ with
\[
 \sum_{a\in A_Q}(X(a)-2-\varepsilon)_+
 \le\rho_\varepsilon H(V)\le\rho_\varepsilon H(X).
 \]
Increase the common threshold until the mixed coefficient is at most
$1-\rho_\varepsilon$, and add. This proves \eqref{ss:shared-hinge}.
Both source classes have used the same target hinge, rather than
independent primal covers.

Finally, the clipped raw boxed dual identity \cite{Synthesis} is
\[
 \mathcal L_{\rm raw}(A,I)=
 \max_{0\le w_p\le1}
 \left(\sum_{a\in A}X(a)-\sum_{b\in I}(X(b)-1)_+\right).
 \]
The weights are on the source primes. Using
$X(a)\le2+\varepsilon+(X(a)-2-\varepsilon)_+$ in this identity
proves the stated fractional bound.
\end{proof}

\subsection{Every fixed positive power cutoff}

\begin{corollary}
\label{ss:fixed-alpha}
Fix $\alpha>0$ and $Y(M)\ge2$ with $\log Y=o(\sqrt{\log M})$.
For distinct sources $A\subseteq[2,M]$ whose prime factors all satisfy
\[
 p\le Y(M)\quad\hbox{or}\quad p\ge M^\alpha,
\]
one has, uniformly over the source set and interval shift,
\[
 \mathcal L_{\rm raw}(A,I)\le(2+\varepsilon)|A|
 \quad\hbox{for all sufficiently large }M,
 \qquad
 \tau(A,I)\le(2+o(1))|A|\quad(|A|\to\infty).
\]
There is no separate bound on either the number of small-prime
occurrences or the number of large-prime occurrences in a source.
\end{corollary}

\begin{proof}
If $\alpha>1$, no source contains a large prime, and the pure-core
result suffices. Otherwise take $Z=M^\alpha$ and
$t=\lfloor1/\alpha\rfloor$. The size bound on each source ensures
$\Omega_R(d)\le t$.
If $r\in R_0$ occurs in $d(a)$, then $r\le p(a)$ and
$rp(a)\le a\le M$, so
$R_0\subseteq[M^\alpha,\sqrt M]$. Repeated prime occurrences are included
in this argument. For $0<\alpha<1/2$, Mertens' reciprocal estimate yields
\[
 \sum_{r\in R_0}\frac1r
 \le\log\frac1{2\alpha}+o(1).
 \]
For $\alpha>1/2$ the set is empty; at $\alpha=1/2$ at most the endpoint
prime occurs and its contribution is $o(1)$.
Choose a fixed integer $L$ with
\[
 L>4\left(1+\max\{0,\log(1/(2\alpha))\}\right).
 \]
The reciprocal condition then holds for large $M$.
All of \eqref{ss:general-conditions} follow from the assumed cutoff
on $Y$, since $t,L,C_*(t,L)$ are fixed. Theorem~\ref{ss:general}
proves the uniform raw bound.

The class is hereditary at the fixed ambient $M$: deleting sources
preserves their prime ranges. Lemma~\ref{lem:tr-hereditary}
therefore supplies the integer conclusion.
Its uniform large-$M$ threshold applies as $|A|\to\infty$, since
distinctness gives $M\ge|A|+1$.
\end{proof}

\subsection{A quantified cutoff below every fixed positive power}

The multiplicity estimate also allows $t$ to grow.
For a support restriction $p\ge Z$, take
$t=\lfloor\log M/\log Z\rfloor$ in the nonempty case $Z\le M$.
As before $R_0\subseteq[Z,\sqrt M]$.
When $Z\to\infty$, Mertens' reciprocal estimate gives
\[
 \sum_{r\in R_0}\frac1r\le\log t+o(1).
 \]
Indeed, if $Z\le\sqrt M$, the main term is
$\log((\log M)/(2\log Z))\le\log t$; otherwise $R_0$ is empty.
Thus the choice
\begin{equation}
 L=\lceil4\log(2t)+4\rceil
 \label{ss:growing-discount}
\end{equation}
satisfies the reciprocal condition for large $M$.

For $t\to\infty$, put $k=t+L-1$. Bounding the sum defining $C_2$ by
$L[k(t+1)]^{L-1}$ gives
\[
 \begin{split}
 \log C_*(t,L)
 &\le\log(t+L)+\log L
       +(L-1)\{\log(t+L-1)+\log(t+1)\}\\
 &\le8(\log t)^2+O(\log t).
 \end{split}
\]
The constant $8$ uses $L=4\log t+O(1)$ from
\eqref{ss:growing-discount}. The exponents of nonexceptional primes
are fixed by the target integer; assigning additional exponent choices
to them would lose this sharper estimate.

\begin{corollary}[An explicit growing-multiplicity class]
\label{ss:subpolynomial}
Fix $0<a<1/2$ and $c>0$ with $8c^2<1-a$. Put
\[
 Y(M)=\exp((\log M)^a),\qquad
 Z(M)=\exp\left(\frac{\log M}
                        {\exp(c\sqrt{\log\log M})}\right).
 \]
If every prime factor of every distinct source in $A\subseteq[2,M]$
is at most $Y(M)$ or at least $Z(M)$, then the uniform raw and
integer conclusions of Corollary~\ref{ss:fixed-alpha} hold.
\end{corollary}

\begin{proof}
Take $t=\lfloor\log M/\log Z\rfloor$ and $L$ from
\eqref{ss:growing-discount}. Then
$\log t=c\sqrt{\log\log M}+o(1)$, and the preceding bound gives
\[
 C_*(t,L)\le(\log M)^{8c^2+o(1)}.
 \]
Furthermore,
\[
 \frac{\log Y}{\log Z}=(\log M)^{a-1+o(1)},\qquad
 \frac{(\log Y)^2}{L\log Z}
       =(\log M)^{2a-1+o(1)}\longrightarrow0.
 \]
All conditions in \eqref{ss:general-conditions} hold, with
$C_*(t,L)\log Y/\log Z\to0$ following from $8c^2+a-1<0$.
Theorem~\ref{ss:general} proves the raw conclusion.
The support class is hereditary at fixed $M$, so
Lemma~\ref{lem:tr-hereditary} gives the integer conclusion.
\end{proof}

For the concrete choice $a=1/3$ and $c=1/4$,
\[
 C_*(t,L)\frac{\log Y}{\log Z}
 \le(\log M)^{-1/6+o(1)}\to0,\qquad
 \frac{(\log Y)^2}{L\log Z}
       =(\log M)^{-1/3+o(1)}\to0.
 \]
Here
$Z=M^{\exp(-\frac14\sqrt{\log\log M})}=M^{o(1)}$.
The theorem therefore does not require a fixed positive exponent
for the large-prime cutoff. It still imposes a substantial gap
between $Y$ and $Z$. The fixed-$\alpha$ theorem cannot be applied
with an arbitrary data-dependent $\alpha$ tending to zero; such
a choice requires the separate quantitative conditions above.

\section{A uniform exclusion for large weighted means}
\label{mo:section}

The following estimates give a further restriction on a possible raw dual
witness. They apply to arbitrary clipped weights, independently of the
structure of the source set.
Let $M\ge4$, let $0\le w_p\le1$ be supported on primes $p\le M$, and set
\[
 X(a)=\sum_{p\le M}w_pv_p(a),\qquad
 \mu=\sum_{p\le M}\frac{w_p}{p-1},\qquad
 \nu=\sum_{p\le M}\frac{w_p^2p}{(p-1)^2}.
\]
Expectations in this section are over a uniform integer in $\{1,\ldots,M\}$.
Put $J=\lfloor\log_2M\rfloor$ and
\begin{equation}
 \epsilon_M=
 \frac{\pi(M)+(J+1)\pi(\sqrt M)}{M}
       +\frac1{\lfloor\sqrt M\rfloor}.
 \label{mo:error}
\end{equation}

\begin{lemma}[Uniform first and second moments]
\label{mo:moments}
One has
\begin{equation}
 \mu-\epsilon_M\le \mathbb EX\le\mu,\qquad
 \mathbb EX^2\le\mu^2+\nu,\qquad
 \nu\le\min\left\{2\mu,\frac76\mu+\frac{49}{48}\right\}.
 \label{mo:bounds}
\end{equation}
Moreover,
\[
 \epsilon_M\le\frac{\pi(M)}M+\frac{J+3}{\sqrt M}
 =O\left(\frac1{\log M}+\frac{\log M}{\sqrt M}\right)=o(1).
\]
\end{lemma}

\begin{proof}
Let $j_p=\lfloor\log M/\log p\rfloor$. Prime-power counting gives
\[
 \mathbb EX=\frac1M\sum_{p\le M}w_p
                    \sum_{j=1}^{j_p}\left\lfloor\frac M{p^j}\right\rfloor.
\]
Compare this with $\mu=\sum_pw_p\sum_{j\ge1}p^{-j}$.
The first-level floor errors total at most $\pi(M)/M$.
At levels $j\ge2$, only primes at most $\sqrt M$ occur, each with
at most $J-1$ further floor errors. For those primes the omitted tail is
\[
 \sum_{j>j_p}p^{-j}
  =\frac1{p^{j_p}(p-1)}
  <\frac p{M(p-1)}\le\frac2M.
\]
For $p>\sqrt M$, the tail starts at level two, and
\[
 \sum_{p>\sqrt M}\frac1{p(p-1)}
 \le\sum_{k>\sqrt M}\frac1{k(k-1)}
 =\frac1{\lfloor\sqrt M\rfloor}.
\]
Adding these bounds proves the mean estimate. The asserted error
asymptotic follows from Chebyshev's upper bound for $\pi$;
see, for example, \cite{GranvilleSound}.

For distinct primes $p,q$, divisibility by both levels $p^j,q^k$ gives
\[
 \mathbb E[v_p(a)v_q(a)]
 \le\sum_{j,k\ge1}\frac1{p^jq^k}
 =\frac1{(p-1)(q-1)}.
\]
At one prime,
\[
 \mathbb E[v_p(a)^2]
 \le\sum_{j\ge1}\frac{2j-1}{p^j}
 =\frac{p+1}{(p-1)^2}.
\]
Expanding $X^2$ proves $\mathbb EX^2\le\mu^2+\nu$.
No independence assumption is used.
The inequality $\nu\le2\mu$ follows from $w_p^2\le w_p$ and
$p/(p-1)\le2$.

For the sharper envelope, every $p\ge7$ contributes at most
$(7/6)w_p/(p-1)$ to $\nu$. The maximum possible excesses at $p=2,3,5$
are respectively $5/6$, $1/6$, and $1/48$. Their sum is $49/48$,
proving the last part of \eqref{mo:bounds}.
\end{proof}

\begin{theorem}[A sufficient mean threshold]
\label{mo:threshold}
For all $M$ with $\epsilon_M\le1/2$, the condition
\begin{equation}
 \mu\ge\frac{241}{136}+4\epsilon_M
 \label{mo:threshold-simple}
\end{equation}
implies, for every positive $M$-term interval $I$,
\begin{equation}
 \sum_{a=1}^M(X(a)-2)_+
 \le\sum_{b\in I}(X(b)-1)_+.
 \label{mo:h2}
\end{equation}
The limiting sufficient threshold is
$241/136=1.7720588235\ldots$.
\end{theorem}

\begin{proof}
For any nonnegative random variable with mean $m$ and second moment $q$,
\[
 \mathbb E\min(X,2)
 =\frac{m+2-\mathbb E|X-2|}{2}
 \ge\frac{m+2-\sqrt{q-4m+4}}2.
\]
Thus $q\le m^2+4m-4$ is sufficient for capped mean at least one.
Put $\epsilon=\epsilon_M$. Since $m\ge\mu-\epsilon$, it is enough that
\[
 \nu\le4\mu-4-(2\mu+4)\epsilon+\epsilon^2.
\]
The sharper envelope in Lemma~\ref{mo:moments} reduces this to
\begin{equation}
 \left(\frac{17}{6}-2\epsilon\right)\mu
 \ge\frac{241}{48}+4\epsilon-\epsilon^2.
 \label{mo:threshold-exact}
\end{equation}
At $\mu=241/136+4\epsilon$, the left side minus the right side equals
\[
 \epsilon\left(\frac{773}{204}-7\epsilon\right)\ge0
 \qquad(0\le\epsilon\le1/2).
\]
The coefficient of $\mu$ is positive, so larger $\mu$ also satisfy
\eqref{mo:threshold-exact}. In this range $\mu-\epsilon>0$, which
justifies using monotonicity of $m^2+4m$ on the nonnegative half-line.

Finally, every prime-power count in $I$ is at least its prefix count,
and hence $\sum_I X\ge\sum_{a\le M}X(a)$. Therefore
\[
 \begin{split}
 &\sum_{a\le M}(X(a)-2)_+-\sum_{b\in I}(X(b)-1)_+\\
 &\qquad\le M-\sum_{a\le M}\min(X(a),2)\le0.
 \end{split}
\]
This proves \eqref{mo:h2}.
\end{proof}

The exact sufficient threshold in \eqref{mo:threshold-exact} is
\[
 \mu\ge
 \frac{241/48+4\epsilon_M-\epsilon_M^2}{17/6-2\epsilon_M}.
\]
The simpler estimate $\mu\ge2+4\epsilon_M$ also follows directly
from $\nu\le2\mu$ and
\[
 \min(x,2)\ge ax-\frac{a^2x^2}{8}\quad(0\le a\le1),
 \qquad a=\frac4{\mu+2}.
\]
These thresholds are sufficient; optimality is not asserted.
Primes greater than $M$ must be omitted when computing $\mu$.
They have zero prefix contribution, and restoring them on the target
side only strengthens \eqref{mo:h2}.

For a distinct source subset $A\subseteq[2,M]$, the hinge inequality
implies that the corresponding raw dual objective is at most $2|A|$.
Thus a dual witness exceeding this bound must lie below the mean
threshold. This does not bound the number of substantial coordinates:
for any integer $N\ge2$, choose $N$ distinct primes greater than $N^2$
and assign each weight one. Their mean is at most $N/(N^2-1)\to0$,
despite $N$ substantial weights. This is an example of weight vectors,
not a counterexample instance. More generally, a mean bound controls only
\[
 \sum_{\{p:w_p\ge\eta\}}\frac1{p-1}\le\frac\mu\eta
 \qquad(\eta>0),
\]
and does not resolve the interaction of a growing core with the other
prime coordinates.

\section{The exact remaining asymptotic target}
\label{ct:remaining}

\subsection{Minimality already removes the rounding problem}

Fix \(c>2\), an ambient \(M\), and an interval \(I\). If an
instance fails \(\tau(A,I)\le c|A|\), replace \(A\) by a failing
subset of smallest cardinality, set \(m=|A|\), and write
\[
 K=\lfloor c(m-1)\rfloor+1.
\]
Every deletion has an integer cover of size at most \(K-1\).
Adding two unused positions would give a cover within the budget
\(\lfloor cm\rfloor\), so the forced-bank arguments apply.
For completeness, their short proofs are retained here.

For every prime occurrence \(q\mid a\), all interval multiples
of \(a/q\) belong to every deletion cover of \(A\setminus\{a\}\).
Otherwise a fresh multiple supplies \(a/q\), and at most one extra
unused position repairs the missing \(q\)-unit. Full supply
\eqref{ct:supply} guarantees that repair. Since there are at least
\(q\) such multiples, all source primes satisfy
\begin{equation}
 p<K<cm.
 \label{ct:prime-cap}
\end{equation}

If \(pq\mid a\), counting occurrences and allowing \(p=q\),
and \(pq\ge K\), then
\begin{equation}
 D_p>v_p\left(\prod_{\substack{b\in I\\a/q\mid b}}b\right)
 \ge qv_p(a/q).
 \label{ct:pair-bank}
\end{equation}
If the first strict inequality failed, the forced bank would
already pay the full \(p\)-demand. There are at least \(pq\ge K\)
targets divisible by \(a/(pq)\); one is fresh relative to a cover
of size \(K-1\). Add it, then repair at most one \(q\)-unit.
This again contradicts minimal failure.

Define the demand core and its complement by
\begin{equation}
 E=\{p\in P:p\le\sqrt K\}\cup\{p\in P:D_p>p\},
 \qquad R=P\setminus E.
 \label{ct:demand-core}
\end{equation}
Every source has at most one \emph{occurrence} from \(R\), and
that prime is its unique largest prime. Indeed, if \(r\in R\)
and another occurrence \(q\ge r\) were in the same source, then
\(rq>K\), while \(D_r\le r\le qv_r(a/q)\), contrary to
\eqref{ct:pair-bank}. In particular,
\begin{equation}
 R\subset(\sqrt K,K),\qquad
 |P|\le\pi(cm)=O_c(m/\log m).
 \label{ct:round-small}
\end{equation}

The quadratic gate forces \(M<16m^2\). For
\(H=\{p:D_p>p\}\), the logarithmic demand budget is
\[
 \sum_{p\in H}p\log p
 <\sum_pD_p\log p
 =\sum_{a\in A}\log a
 \le m\log M.
\]
Chebyshev's bound gives \(p_j\gg j\log(j+1)\) for the
\(j\)-th prime. Thus the left side, for \(h=|H|\), is at least
a constant times \(h^2\log^2(h+1)\). Separating
\(h<m^{1/4}\) from its complement gives
\begin{equation}
 |E|=O_c\bigl(\sqrt{m/\log m}\bigr).
 \label{ct:core-size}
\end{equation}
This reproduces the sharper structural bound in \cite{Synthesis}
without an additive-budget parameter. Minimality is always taken
at fixed \(M\); no interval is shortened after a deletion.

\subsection{A sufficient inequality for full asymptotic closure}

Fix \(\delta>0\) and take \(c=2+3\delta\) in the preceding
minimality argument. The remaining analytic goal is the following
uniform estimate for the just-described minimal-obstruction class:
\begin{equation}
 \boxed{\quad
 \sup_{0\le w_p\le1}\Delta_{2+\delta}(A,I,w)=o(m).
 \quad}
 \label{ct:open-hinge}
\end{equation}
The supremum is over \emph{all} clipped core weights, with no
bound on how many exceed a small threshold. The estimate must be
uniform in \(M<16m^2\), the interval shift, the distinct sources,
and their prime arithmetic.

To see that this goal suffices, suppose arbitrarily large instances
failed the budget \((2+3\delta)n\), and pass at fixed ambient
length to a minimal failing subset of size \(m\). The quadratic
gate gives \(m>\sqrt M/4\), so \(m\to\infty\). By
\eqref{ct:open-hinge} and duality,
\(L(A,I)\le(2+\delta)m+o(m)\). By
\eqref{ct:round-small} and \eqref{ct:round}, the same bound holds
for \(\tau(A,I)\), contradicting the chosen budget for large
\(m\). Since \(\delta\) is arbitrary, the unrestricted
\((2+o(1))n\) assertion would follow.

The new controlled-core theorem bounds the supremum restricted to
vectors having at most
\[
 k\le \frac{(1-\gamma)\log\log m}{\log D_\delta},
 \qquad 0<\gamma<1,
\]
core weights exceeding \(c_\delta/(k+1)\) by \(o(m)\). Arbitrarily many
substantial weights in \(R\) are allowed. The unrestricted demand
core in \eqref{ct:core-size} can be much larger, so the result
does not establish \eqref{ct:open-hinge} for every vector.

The separate moment theorem excludes dual vectors with mean
\(\sum_{p\le M}w_p/(p-1)\ge241/136+o(1)\).
The separated-support theorem excludes its stated source classes.
Combining these exclusions does not show that every remaining
instance belongs to one of them.

The direct construction and its strengthening in
Corollary~\ref{cor:sc-chebyshev-core} give a further source-level
restriction, independent of the weights. For the sufficiently
small absolute \(c_0>0\) of that corollary, the demand core of a
sufficiently large minimal asymptotic obstruction must satisfy
\[
 |E|>c_0\frac{(\log M)^2}{\log\log M}.
\]
Indeed, the one-outside-prime property already holds by
\eqref{ct:demand-core}, so a smaller core would place that source
set inside the new integer theorem. Consequently any hypothetical
obstruction sequence has
\[
 \liminf\frac{\log|E|}{\log\log M}\ge2,
 \qquad |E|=O_c\bigl(\sqrt{m/\log m}\bigr).
\]
There is still a large gap between these necessary bounds. This is
not a proof that such an obstruction exists.

\subsection{What has been closed and what still has to change}

The following steps are now established within this proof programme:

\begin{enumerate}
\item The diffuse core and the one-prime remainder share a single
target budget. No separate core cover is added to a remainder cover.
\item All raw prime-power tails are retained. Capping is an exact
intermediate device, followed by a justified simultaneous restoration.
\item A controlled substantial core incurs a finite upper-set error,
independent of \(M\). Each remainder prime pays that error once,
regardless of its exponent range.
\item Fixed-ambient minimality reduces integer rounding to \(o(m)\).
The raw-to-integer passage is no longer an unproved final step.
\item A direct common-core construction handles the core size
bound \(c_0(\log M)^2/\log\log M\), while retaining arbitrary
outside prime powers. Nested matching supplies all outside demands;
exact core-part bins and a sparse exceptional repair supply the
shared core demands.
\item In the separated-support regime, arbitrary pure-core sources
are absorbed, positive certificate overlap is bounded, and the
old positive lower bound on the fixed large-prime exponent is removed.
\end{enumerate}

For the unrestricted asymptotic target, the next proof must control
the joint contribution of many substantial core coordinates more
economically than the variation bound
\(B_{\delta,k}(1+|R|)\). It is sufficient to prove
\eqref{ct:open-hinge}; it is not necessary to bound each upper-set
discrepancy separately. Any such proof must preserve the shared
target budget and the dependence between the core, the remainders,
and the source demands.

For the exact target \(2n\), even a proof of
\eqref{ct:open-hinge} would leave an additional task: removing the
\(o(n)\) slack and the rounding loss, or constructing an exact
integer cover directly. The present note therefore distinguishes
the exact and asymptotic endpoints throughout.

\subsection{Two shortcuts that do not establish the missing step}

The upper set of smooth parts cannot always be counted in a prefix
as though that prefix minimized its count among intervals of equal
length. The set of integers divisible by \(2\) or \(3\) contains
two members of \(\{1,2,3\}\), but only one member of
\(\{5,6,7\}\). This disproves that intermediate count claim;
it is not a counterexample to problem 708. The finite-difference
error in the controlled-core proof is what repairs this step.

Also, the statement ``for every fixed \(\alpha>0\)'' in the
separated-support theorem cannot be used with an arbitrary
instance-dependent \(\alpha\to0\). The shrinking-exponent example
is valid because its quantitative error conditions are checked
separately. Primes between its stated bounds \(Y(M)\) and \(Z(M)\)
remain outside that source class.

The archive contains the complete proof notes, their mathematical
audit records, and the inherited inputs. No finite search is used
as evidence of a universal conclusion in this continuation.


\clearpage
\begin{thebibliography}{9}
\bibitem{Wang}
Z. Wang, \emph{Fractional-to-Integer Transference and Five-Prime
Bounds for Interval Divisor Covers}, public working manuscript,
version 1.0.0, 21 October 2026. Appendix A gives the quadratic
coefficient-two gate.
\href{https://github.com/science-prince/erdos-708-interval-divisor-covers}{Primary repository}
and \href{https://github.com/science-prince/erdos-708-interval-divisor-covers/blob/main/docs/RESULTS.md}{result catalogue}.
Consulted 23 November 2026.

\bibitem{Synthesis}
\emph{Erd\H{o}s Problem 708: Synthesis and Further Partial Results},
AI-assisted working note prepared for Denis Paliy, 23 November 2026,
27 pp. Included as \texttt{inherited/erdos708\_synthesis.pdf} in the
companion archive, together with the analytic proof notes on which
the present continuation depends. Internal research document;
no external review or priority claim.

\bibitem{Lagarias}
J. C. Lagarias, \emph{Math 675: Combinatorial Sieve: Fundamental
Lemma}, theorem and interval example, pp. 1--2.
\href{https://websites.umich.edu/~lagarias/675notes/Estimate5-2015.pdf}{University of Michigan lecture notes}.

\bibitem{GranvilleSound}
A. Granville and K. Soundararajan, \emph{Multiplicative number
theory: The pretentious approach}, Chapter 1, Theorem 1.4.2.
\href{https://dms.umontreal.ca/~andrew/PDF/Book.To1.4.pdf}{Author-hosted text}.

\bibitem{TaoMertens}
T. Tao, \emph{Mertens' theorems}, 11 December 2013.
\href{https://terrytao.wordpress.com/2013/12/11/mertens-theorems/}{Author's exposition}.
\end{thebibliography}
\end{document}
